%%%%%%%%%%%%%%%%%%%%%%%%%%%%%%%%%%%%%%%%%%%%%%%%%%%%%%%%%%%%%%%%%%%%%%%%%%%%%%%
% Harvard Academic Notes - 통합 마스터 템플릿
% 모든 강의 노트에 적용되는 통일된 스타일
% 버전: 2.1 - 가독성 개선 (선택적 최적화)
% 최종 수정일: 2025-11-17
%%%%%%%%%%%%%%%%%%%%%%%%%%%%%%%%%%%%%%%%%%%%%%%%%%%%%%%%%%%%%%%%%%%%%%%%%%%%%%%

\documentclass[11pt,a4paper]{article}

%========================================================================================
% 기본 패키지
%========================================================================================

% --- 한국어 지원 ---
\usepackage{kotex}

% --- 페이지 레이아웃 ---
\usepackage[top=20mm, bottom=20mm, left=20mm, right=18mm]{geometry}
\usepackage{setspace}
\onehalfspacing                      % 1.5배 줄간격
\setlength{\parskip}{0.5em}          % 문단 간격
\setlength{\parindent}{0pt}          % 들여쓰기 없음

% --- 표 관련 ---
\usepackage{booktabs}              % 고품질 표
\usepackage{tabularx}              % 자동 너비 조절 표
\usepackage{array}                 % 표 컬럼 확장
\usepackage{longtable}             % 여러 페이지 표
\renewcommand{\arraystretch}{1.1}  % 표 행간 조절

%========================================================================================
% 헤더 및 푸터
%========================================================================================

\usepackage{fancyhdr}
\pagestyle{fancy}
\fancyhf{}
\fancyhead[L]{\small\textit{FIN-571: Financial Institutions & Monetary Policy}}
\fancyhead[R]{\small\textit{Module 05}}
\fancyfoot[C]{\thepage}
\renewcommand{\headrulewidth}{0.5pt}
\renewcommand{\footrulewidth}{0.3pt}

% 첫 페이지는 헤더 없음
\fancypagestyle{firstpage}{
    \fancyhf{}
    \fancyfoot[C]{\thepage}
    \renewcommand{\headrulewidth}{0pt}
}

%========================================================================================
% 색상 정의 (파스텔 톤 + 다크모드 호환)
%========================================================================================

\usepackage[dvipsnames]{xcolor}

% 밝은 배경용 파스텔 색상
\definecolor{lightblue}{RGB}{220, 235, 255}      % 부드러운 파랑
\definecolor{lightgreen}{RGB}{220, 255, 235}     % 부드러운 초록
\definecolor{lightyellow}{RGB}{255, 250, 220}    % 부드러운 노랑
\definecolor{lightpurple}{RGB}{240, 230, 255}    % 부드러운 보라
\definecolor{lightgray}{gray}{0.95}              % 밝은 회색
\definecolor{lightpink}{RGB}{255, 235, 245}      % 부드러운 핑크
\definecolor{boxgray}{gray}{0.95}
\definecolor{boxblue}{rgb}{0.9, 0.95, 1.0}
\definecolor{boxred}{rgb}{1.0, 0.95, 0.95}

% 진한 색상 (테두리/제목용)
\definecolor{darkblue}{RGB}{50, 80, 150}
\definecolor{darkgreen}{RGB}{40, 120, 70}
\definecolor{darkorange}{RGB}{200, 100, 30}
\definecolor{darkpurple}{RGB}{100, 60, 150}

%========================================================================================
% 박스 환경 (tcolorbox) - 6가지 타입
%========================================================================================

\usepackage[most]{tcolorbox}
\tcbuselibrary{skins, breakable}

% 1. 개요 박스 (강의 시작 부분)
\newtcolorbox{overviewbox}[1][]{
    enhanced,
    colback=lightpurple,
    colframe=darkpurple,
    fonttitle=\bfseries\large,
    title=📚 강의 개요,
    arc=3mm,
    boxrule=1pt,
    left=8pt,
    right=8pt,
    top=8pt,
    bottom=8pt,
    breakable,
    #1
}

% 2. 요약 박스
\newtcolorbox{summarybox}[1][]{
    enhanced,
    colback=lightblue,
    colframe=darkblue,
    fonttitle=\bfseries,
    title=📝 핵심 요약,
    arc=2mm,
    boxrule=0.7pt,
    left=6pt,
    right=6pt,
    top=6pt,
    bottom=6pt,
    breakable,
    #1
}

% 3. 핵심 정보 박스
\newtcolorbox{infobox}[1][]{
    enhanced,
    colback=lightgreen,
    colframe=darkgreen,
    fonttitle=\bfseries,
    title=💡 핵심 정보,
    arc=2mm,
    boxrule=0.7pt,
    left=6pt,
    right=6pt,
    top=6pt,
    bottom=6pt,
    breakable,
    #1
}

% 4. 주의사항 박스
\newtcolorbox{warningbox}[1][]{
    enhanced,
    colback=lightyellow,
    colframe=darkorange,
    fonttitle=\bfseries,
    title=⚠️ 주의사항,
    arc=2mm,
    boxrule=0.7pt,
    left=6pt,
    right=6pt,
    top=6pt,
    bottom=6pt,
    breakable,
    #1
}

% 5. 예제 박스
\newtcolorbox{examplebox}[1][]{
    enhanced,
    colback=lightgray,
    colframe=black!60,
    fonttitle=\bfseries,
    title=📖 예제,
    arc=2mm,
    boxrule=0.7pt,
    left=6pt,
    right=6pt,
    top=6pt,
    bottom=6pt,
    breakable,
    #1
}

% 6. 정의 박스
\newtcolorbox{definitionbox}[1][]{
    enhanced,
    colback=lightpink,
    colframe=purple!70!black,
    fonttitle=\bfseries,
    title=📌 정의,
    arc=2mm,
    boxrule=0.7pt,
    left=6pt,
    right=6pt,
    top=6pt,
    bottom=6pt,
    breakable,
    #1
}

% 7. 중요 박스 (importantbox - warningbox와 유사)
\newtcolorbox{importantbox}[1][]{
    enhanced,
    colback=boxred,
    colframe=red!70!black,
    fonttitle=\bfseries,
    title=⚠️ 매우 중요,
    arc=2mm,
    boxrule=0.7pt,
    left=6pt,
    right=6pt,
    top=6pt,
    bottom=6pt,
    breakable,
    #1
}

% 8. cautionbox (warningbox와 동일)
\let\cautionbox\warningbox
\let\endcautionbox\endwarningbox

%========================================================================================
% 코드 블록 설정 (밝은 배경)
%========================================================================================

\usepackage{listings}

\definecolor{codegray}{rgb}{0.5,0.5,0.5}
\definecolor{codepurple}{rgb}{0.58,0,0.82}
\definecolor{backcolour}{rgb}{0.95,0.95,0.95}

\lstset{
    basicstyle=\ttfamily\small,
    backgroundcolor=\color{lightgray},
    keywordstyle=\color{darkblue}\bfseries,
    commentstyle=\color{darkgreen}\itshape,
    stringstyle=\color{purple!80!black},
    numberstyle=\tiny\color{black!60},
    numbers=left,
    numbersep=8pt,
    breaklines=true,
    breakatwhitespace=false,
    frame=single,
    frameround=tttt,
    rulecolor=\color{black!30},
    captionpos=b,
    showstringspaces=false,
    tabsize=2,
    xleftmargin=15pt,
    xrightmargin=5pt,
    escapeinside={\%*}{*)}
}

% Python 코드 스타일
\lstdefinestyle{pythonstyle}{
    language=Python,
    morekeywords={self, True, False, None},
}

% SQL 코드 스타일
\lstdefinestyle{sqlstyle}{
    language=SQL,
    morekeywords={SELECT, FROM, WHERE, JOIN, GROUP, BY, ORDER, HAVING},
}

%========================================================================================
% 목차 스타일링
%========================================================================================

\usepackage{tocloft}
\renewcommand{\cftsecleader}{\cftdotfill{\cftdotsep}}
\setlength{\cftbeforesecskip}{0.4em}
\renewcommand{\cftsecfont}{\bfseries}
\renewcommand{\cftsubsecfont}{\normalfont}

%========================================================================================
% 표 및 그림
%========================================================================================

\usepackage{graphicx}              % 이미지
\usepackage{adjustbox}             % 표/박스 크기 조절

% 표 캡션 스타일
\usepackage{caption}
\captionsetup[table]{
    labelfont=bf,
    textfont=it,
    skip=5pt
}
\captionsetup[figure]{
    labelfont=bf,
    textfont=it,
    skip=5pt
}

%========================================================================================
% 수학
%========================================================================================

\usepackage{amsmath, amssymb, amsthm}

% 정리 환경
\theoremstyle{definition}
\newtheorem{theorem}{정리}[section]
\newtheorem{lemma}[theorem]{보조정리}
\newtheorem{proposition}[theorem]{명제}
\newtheorem{corollary}[theorem]{따름정리}
\newtheorem{definition}{정의}[section]
\newtheorem{example}{예제}[section]

%========================================================================================
% 하이퍼링크
%========================================================================================

\usepackage[
    colorlinks=true,
    linkcolor=blue!80!black,
    urlcolor=blue!80!black,
    citecolor=green!60!black,
    bookmarks=true,
    bookmarksnumbered=true,
    pdfborder={0 0 0}
]{hyperref}

% PDF 메타데이터는 각 문서에서 설정
\hypersetup{
    pdftitle={FIN-571: Financial Institutions & Monetary Policy - Module 05},
    pdfauthor={강의 노트},
    pdfsubject={Academic Notes}
}

%========================================================================================
% 기타 유용한 패키지
%========================================================================================

\usepackage{enumitem}              % 리스트 커스터마이징
\setlist{nosep, leftmargin=*, itemsep=0.3em}

\usepackage{microtype}             % 타이포그래피 개선
\usepackage{footnote}              % 각주 개선
\usepackage{url}                   % URL 줄바꿈
\urlstyle{same}

%========================================================================================
% 사용자 정의 명령어
%========================================================================================

% 강조 텍스트
\newcommand{\important}[1]{\textbf{\textcolor{red!70!black}{#1}}}
\newcommand{\keyword}[1]{\textbf{#1}}
\newcommand{\term}[1]{\textit{#1}}
\newcommand{\code}[1]{\texttt{#1}}

% 용어 설명 (인라인)
\newcommand{\defterm}[2]{\textbf{#1}\footnote{#2}}

% 섹션 시작 전 페이지 분리
\newcommand{\newsection}[1]{\newpage\section{#1}}

%========================================================================================
% 문서 제목 스타일
%========================================================================================

\usepackage{titling}
\pretitle{\begin{center}\LARGE\bfseries}
\posttitle{\par\end{center}\vskip 0.5em}
\preauthor{\begin{center}\large}
\postauthor{\end{center}}
\predate{\begin{center}\large}
\postdate{\par\end{center}}

%========================================================================================
% 섹션 제목 간격
%========================================================================================

\usepackage{titlesec}
\titlespacing*{\section}{0pt}{1.5em}{0.8em}
\titlespacing*{\subsection}{0pt}{1.2em}{0.6em}
\titlespacing*{\subsubsection}{0pt}{1em}{0.5em}

%========================================================================================
% 메타 정보 박스 명령어
%========================================================================================

\newcommand{\metainfo}[4]{
\begin{tcolorbox}[
    colback=lightpurple,
    colframe=darkpurple,
    boxrule=1pt,
    arc=2mm,
    left=10pt,
    right=10pt,
    top=8pt,
    bottom=8pt
]
\begin{tabular}{@{}rl@{}}
▣ \textbf{강의명:} & #1 \\[0.3em]
▣ \textbf{주차:} & #2 \\[0.3em]
▣ \textbf{교수명:} & #3 \\[0.3em]
▣ \textbf{목적:} & \begin{minipage}[t]{0.75\textwidth}#4\end{minipage}
\end{tabular}
\end{tcolorbox}
}

%========================================================================================
% 끝
%========================================================================================


\begin{document}

\thispagestyle{firstpage}

\metainfo{금융 시장과 통화 정책}{Module 1: 금리 및 증권화}{Prof. Ralf Meisenzahl}{중앙은행과 통화정책의 핵심 개념, 금리의 종류와 작동 방식, 주요 금융 시장 및 금융 상품에 대한 심층적인 이해를 목표로 함.}

\tableofcontents

\newpage

\section{개요}
이 문서는 중앙은행과 통화정책의 핵심 개념을 다루는 모듈 1의 내용을 정리한 것입니다.
통화정책이 어떻게 실행되고 거시 경제에 영향을 미치는지 이해하기 위한 기초 지식을 제공합니다.
특히, 금리의 종류와 작동 방식, 주요 금융 시장 및 금융 상품에 대한 심층적인 이해를 목표로 합니다.
단기 및 장기 금리, 국채, 환매조건부채권(Repo), 머니 마켓 뮤추얼 펀드 등 다양한 시장과 도구를 학습하여,
통화정책의 복잡한 구현 과정을 파악하고 경제 예측에 활용되는 지표들을 분석할 수 있습니다.
이 모듈은 다음 모듈에서 다룰 통화정책의 구체적인 도구와 구현 방식을 이해하기 위한 필수적인 배경 지식을 제공합니다.

\newpage
\section{용어 정리}
다음 표는 본 문서에서 다루는 주요 용어들을 쉽게 이해할 수 있도록 정리한 것입니다.

\begin{table}[h!]
\centering
\caption{주요 금융 용어 정리}
\label{tab:terms}
\begin{adjustbox}{width=\textwidth,center}
\begin{tabular}{lllc}
\toprule
\textbf{용어 (원어)} & \textbf{쉬운 설명} & \textbf{비고} \\
\midrule
연방기금 시장 (Federal Funds Market) & 은행들이 연방준비제도에 예치된 초과 지급준비금을 서로 빌려주고 빌리는 시장. & 주로 익일물(overnight) 대출. \\
연방기금 금리 (Federal Funds Rate) & 연방기금 시장에서 은행 간에 초과 지급준비금을 거래할 때 적용되는 금리. & 중앙은행의 주요 정책 목표 금리. \\
환매조건부채권 (Repo) & 채권을 담보로 현금을 빌리고, 다음 날 채권을 다시 사들이기로 약속하는 거래. & 담보가 있는 익일물 대출 시장. \\
리보 (LIBOR) & 런던 은행 간 대출 금리. 은행 간 무담보 대출에 적용되는 금리. & 과거 조작 논란으로 SOFR로 대체 중. \\
소퍼 (SOFR) & 담보부 익일물 자금 조달 금리. 국채를 담보로 하는 환매조건부채권 시장 금리. & LIBOR를 대체하는 새로운 무위험 기준 금리. \\
명목 금리 (Nominal Interest Rate) & 인플레이션을 고려하지 않은, 표면적인 금리. & 은행 예금 이자율 등. \\
실질 금리 (Real Interest Rate) & 인플레이션을 고려하여 구매력 변화를 반영한 금리. & 실제 투자 수익률을 나타냄. \\
피셔 방정식 (Fisher Equation) & 명목 금리, 실질 금리, 예상 인플레이션 간의 관계를 설명하는 방정식. & 명목 금리 $\approx$ 실질 금리 + 예상 인플레이션. \\
물가연동국채 (TIPS) & 원금 가치가 인플레이션에 따라 조정되는 국채. & 실질 금리를 지급하며 인플레이션 헤지 기능. \\
국채 (Treasury Securities) & 미국 재무부가 발행하는 채권. 신용 위험이 없어 무위험 자산으로 간주. & T-Bill, T-Note, T-Bond로 구분. \\
무위험 금리 (Risk-Free Rate) & 신용 위험이 없는 투자에 대한 이론적인 수익률. & 보통 미국 국채 금리를 기준으로 삼음. \\
금리 기간 구조 (Term Structure of Interest Rates) & 동일 발행 주체의 만기별 금리 비교. & 수익률 곡선(Yield Curve)으로 시각화. \\
수익률 곡선 (Yield Curve) & 만기가 다른 채권들의 수익률을 연결한 곡선. & 경제 상황 및 미래 금리 기대를 반영. \\
머니 마켓 뮤추얼 펀드 (Money Market Mutual Funds) & 초단기, 초안전 자산에 투자하는 뮤추얼 펀드. & 은행 예금의 대안으로 활용. \\
지급준비금 (Reserves) & 은행이 예금 인출 등에 대비하여 중앙은행에 예치하는 현금. & 법정 지급준비금과 초과 지급준비금으로 나뉠 수 있음. \\
법정 지급준비금 (Required Reserves) & 법으로 정해진 최소한의 지급준비금. & 중앙은행이 결정. \\
초과 지급준비금 (Excess Reserves) & 법정 지급준비금을 초과하여 보유하는 지급준비금. & 은행의 유동성 수요에 대비. \\
헤어컷 (Haircut) & 담보 가치보다 적은 금액을 대출해주는 비율. & 담보 가치 하락 위험을 줄이기 위함. \\
\bottomrule
\end{tabular}
\end{adjustbox}
\end{table}

\newpage
\section{모듈 1: 금리 (Interest Rates)}

\subsection{개요}
이 모듈은 중앙은행, 통화정책, 그리고 거시 경제에 대한 이해를 돕기 위한 첫걸음입니다. 우리는 통화정책이 어떻게 실행되고, 중앙은행의 결정이 궁극적으로 거시 경제에 어떻게 전달되는지 논의할 것입니다. 이 과정을 통해 통화정책의 목표에 대한 깊이 있는 이해를 얻게 될 것입니다.

\begin{infobox}[title=\textbf{통화정책의 핵심 질문}]
중앙은행이 "금리를 25bp 인상한다"거나 "장기 금리에 하방 압력을 가하기 위해 자산 매입 프로그램을 확대한다"는 발표를 들었을 때, 이것이 실제로 무엇을 의미할까요? 중앙은행은 이러한 정책을 어떻게 실행하며, 어떤 금리가 왜 영향을 받을까요?
\end{infobox}

통화정책의 실행은 여러 금융 시장과 금융 상품을 포함하는 복잡한 주제입니다. 이 모듈은 통화정책이 어떤 금리를 목표로 하는지 정확히 설명하는 데 필요한 배경 지식을 제공합니다.

\begin{warningbox}[title=\textbf{주의: 모듈 1의 특성}]
이 모듈의 영상들은 다음 모듈에서 통화정책을 본격적으로 다루기 때문에 다소 단절된 느낌을 줄 수 있습니다. 각 영상이 전체 그림에 어떻게 들어맞는지 궁금할 수 있습니다. 하지만 이 영상들은 통화정책과 그 실행을 이해하기 위한 배경 지식을 제공하는 것임을 기억하세요. 나중에 특정 통화정책 도구나 금융 위기를 논의할 때 언제든지 이 영상들로 돌아와 개념을 다시 확인할 수 있습니다.
\end{warningbox}

\subsubsection*{모듈 1 영상 개요}
\begin{itemize}
    \item \textbf{단기 자금 시장:} 첫 두 영상은 연방기금 시장(Federal Funds Market)과 환매조건부채권(Repo) 시장이라는 두 가지 단기 자금 시장을 소개합니다. 이 시장들은 금융 기관들이 현금을 빌리고 빌려주는 익일물(overnight) 시장이며, 통화정책의 주요 목표 대상입니다.
    \item \textbf{연방준비제도의 목표:} 연방준비제도(Federal Reserve)는 연방기금 시장의 금리인 연방기금 금리(Federal Funds Rate)를 목표 범위로 설정하여 관리합니다. 환매조건부채권 시장은 연방기금 시장의 대안이므로, 연방준비제도는 모든 익일물 금리가 통화정책 목표와 일치하도록 환매조건부채권 시장의 금리도 목표로 합니다.
    \item \textbf{런던 은행 간 대출 금리 (LIBOR):} 다음 영상에서는 또 다른 익일물 금리인 런던 은행 간 대출 금리(London Interbank Offered Rate, LIBOR)를 환매조건부채권 시장 금리와 비교합니다. 이 금리들은 기업 대출의 벤치마크가 됩니다.
    \item \textbf{주요 금융 상품:} 통화정책의 실행과 전달에 핵심적인 세 가지 금융 상품을 살펴봅니다.
    \begin{enumerate}
        \item \textbf{정부 증권 (Government Securities):} 중앙은행은 통화정책을 실행하기 위해 정부 증권을 사고팔기 때문에, 이 증권이 무엇이고 어떻게 작동하는지 이해하는 것이 중요합니다. 미국 정부 증권은 채무 불이행 위험이 없으므로, 금리의 구성 요소(예: 인플레이션 반영분)를 분석할 수 있습니다.
        \item \textbf{금리 기간 구조 (Term Structure of Interest Rates):} 금리가 인플레이션과 미래에 대한 기대를 반영하므로, 만기가 짧은 정부 증권과 만기가 긴 정부 증권의 금리를 비교하여 시장의 미래 경제 발전 기대를 이해할 수 있습니다. 이는 중앙은행, 투자 은행, 금융 분석가들이 면밀히 주시하는 지표입니다.
        \item \textbf{머니 마켓 뮤추얼 펀드 (Money Market Mutual Funds):} 머니 펀드는 환매조건부채권 시장에서 가장 큰 익일물 현금 대출 기관입니다. 2008년 금융 위기 이후 연방준비제도는 이러한 펀드를 통해 통화정책을 수행하기 시작했습니다.
        \item \textbf{증권화 (Securitization):} 증권화는 신용 공급을 증가시켰습니다. 특히 주택담보부증권(Mortgage-Backed Securities, MBS)과 같은 증권화된 모기지 풀은 매우 중요합니다. 통화정책은 신용 공급을 줄이거나 늘림으로써 경제에 영향을 미치고, 연방준비제도는 특정 MBS를 매입할 수 있으므로, 증권화 시장은 통화정책 고려 사항에서 중요한 요소입니다.
    \end{enumerate}
\end{itemize}
요약하자면, 이 첫 번째 모듈에서 논의되는 모든 시장과 상품은 통화정책 및 그 실행과 직접적으로 관련되어 있습니다. 각 영상은 다음 모듈에서 더 명확해질 큰 그림의 작은 조각들을 제공합니다.

\newpage
\section{단기 금리 (Short-term Rates)}

\subsection{연방기금 시장 (The Market for Federal Funds)}
이 강의에서는 연방기금 시장의 배경과 작동 방식을 논의합니다.

\subsubsection*{은행이 지급준비금을 보유하는 이유}
\begin{itemize}
    \item \textbf{은행의 사업 모델:} 은행은 단기로 자금을 빌려(예금 유치) 장기로 대출(대출 실행)하는 사업 모델을 가지고 있습니다.
    \item \textbf{현금 수요 (유동성):} 이러한 사업 모델은 은행에 현금, 즉 유동성(liquidity)에 대한 수요를 발생시킵니다. 예를 들어, 예금은 언제든지 인출될 수 있으므로, 은행은 유동성 수요를 충족시키기 위해 일정량의 현금을 보유해야 합니다. 이 현금은 연방준비제도(Federal Reserve)의 지급준비금 계좌에 예치될 수 있습니다.
\end{itemize}

\subsubsection*{지급준비금의 종류: 법정 지급준비금과 초과 지급준비금}
은행은 얼마나 많은 지급준비금을 보유해야 할까요? 이는 어려운 질문입니다.

\begin{itemize}
    \item \textbf{법정 지급준비금 (Required Reserves):}
    \begin{itemize}
        \item 미국에서는 은행이 연방준비제도에 최소한의 지급준비금을 보유하도록 법적으로 요구됩니다. 이를 법정 지급준비금이라고 합니다.
        \item 이 지급준비금 요건은 은행이 일상적인 현금 수요를 충족시킬 수 있도록 설정되며, 연방준비제도가 결정합니다.
        \item \textbf{예시:} 코로나19 팬데믹 이전에는 예금의 3\%였습니다. 팬데믹 동안 은행을 지원하기 위해 이 요건은 0\%로 설정되었습니다.
    \end{itemize}
    \item \textbf{초과 지급준비금 (Excess Reserves):}
    \begin{itemize}
        \item 은행은 일반적으로 법정 지급준비금보다 더 많은 지급준비금을 보유합니다. 이를 초과 지급준비금이라고 합니다.
        \item 은행은 예상치 못한 유동성 수요를 충족시키기 위해 초과 지급준비금을 보유합니다.
        \item \textbf{예시:} 코로나19 팬데믹 초기에 미국 기업들은 신용 한도에서 약 6천억 달러를 인출했습니다. 은행들은 이 예측 불가능한 사건에 대비할 수 없었지만, 단기간에 고객에게 6천억 달러의 현금을 제공해야 했습니다. 이것이 초과 지급준비금의 필요성을 보여줍니다.
    \end{itemize}
\end{itemize}

\subsubsection*{은행의 지급준비금 수준 선택: 유동성과 기회비용}
은행은 지급준비금 수준을 어떻게 선택할까요?

\begin{itemize}
    \item 은행은 현금을 지급준비금 형태로 보유하거나, 이 현금을 투자할 수 있습니다. 예를 들어, 은행에 이자 수입을 안겨주는 대출에 투자할 수 있습니다.
    \item \textbf{상충 관계 (Trade-offs):}
    \begin{itemize}
        \item 지급준비금을 보유하는 것은 일반적으로 대출에 비해 낮은 수입을 창출합니다.
        \item 그러나 예상치 못한 유동성 수요에 직면하면, 은행은 막판에 자금을 빌려야 할 수 있으며, 이는 비용이 많이 들 수 있습니다.
        \item 따라서 은행은 유동성 수요와 이자 지급 자산에 투자했을 때 얻을 수 있었던 수익(기회비용) 사이에서 상충 관계에 직면합니다.
    \end{itemize}
\end{itemize}

\begin{examplebox}[title=\textbf{초과 지급준비금의 기회비용 예시}]
은행이 100억 원의 초과 지급준비금을 보유하고 있다고 가정해 봅시다. 이 돈을 대출해 주면 연 5\%의 이자를 받을 수 있지만, 지급준비금으로 보유하면 연 0.1\%의 이자만 받습니다. 이 경우, 대출을 통해 얻을 수 있었던 4.9\%의 이자 수익이 초과 지급준비금 보유의 기회비용이 됩니다. 은행은 이 기회비용과 예상치 못한 유동성 부족으로 인한 위험 및 비용을 저울질하여 최적의 지급준비금 수준을 결정합니다.
\end{examplebox}

\subsubsection*{미국 상업은행의 지급준비금 보유 현황}
\begin{itemize}
    \item \textbf{역사적으로 낮은 수준:} 1960년대부터 2007년까지 지급준비금은 거의 증가하지 않았습니다.
    \item \textbf{2008년 금융 위기 이후 급증:}
    \begin{itemize}
        \item 2007년 약 400억 달러에서 2014년 거의 3조 달러로, 2021년 3월에는 3.7조 달러로 크게 증가했습니다.
        \item \textbf{급증의 원인:}
        \begin{enumerate}
            \item \textbf{저금리 환경 (2008-2022):} 대출과 같은 이자 지급 자산의 수익률이 감소하여 초과 지급준비금 보유의 기회비용이 낮아졌습니다.
            \item \textbf{지급준비금에 대한 이자 지급:} 연방준비제도가 지급준비금에 이자를 지급하기 시작하여 초과 지급준비금 보유의 기회비용이 더욱 낮아졌습니다.
            \item \textbf{금융 위기 경험:} 금융 위기 동안 은행들은 대규모 현금 수요를 경험했고, 살아남은 은행들은 안전장치로 더 많은 지급준비금을 보유하기로 결정했습니다.
        \end{enumerate}
    \end{itemize}
\end{itemize}

\subsubsection*{연방기금 시장의 작동 방식}
은행이 유동성 수요를 충족시킬 만큼 충분한 지급준비금을 가지고 있지 않을 때 어떻게 될까요?

\begin{itemize}
    \item 은행은 다른 은행으로부터 지급준비금을 빌릴 수 있습니다. 이 거래가 \textbf{연방기금 시장 (Federal Funds Market)}에서 이루어집니다.
    \item \textbf{특징:}
    \begin{itemize}
        \item 대출은 보통 익일물(overnight)이며, 담보를 요구하지 않습니다.
        \item 이는 은행이 채무 불이행을 하지 않을 것이라는 높은 신뢰를 요구합니다.
        \item 이 시장에서 빌리는 비용을 \textbf{연방기금 금리 (Federal Funds Rate)}라고 합니다.
        \item 연방기금 시장은 익일물 현금을 제공하는 소위 \textbf{머니 마켓 (Money Market)}입니다.
    \end{itemize}
\end{itemize}

\subsubsection*{연방기금 시장의 중요성}
연방기금 시장은 두 가지 핵심 기능을 수행합니다.

\begin{enumerate}
    \item \textbf{은행 간 자금 재분배:}
    \begin{itemize}
        \item 일부 은행은 대출 기회가 많지만 현금이 부족합니다. 다른 은행은 지급준비금 형태로 현금이 많지만 대출 기회가 적습니다.
        \item 연방기금 시장을 통해 후자의 은행은 전자의 은행에 현금을 빌려줄 수 있습니다. 이를 통해 전자의 은행은 대출 수요를 충족시킬 수 있습니다.
    \end{itemize}
    \item \textbf{경제 전반의 신용 비용에 영향:}
    \begin{itemize}
        \item 연방기금 금리는 경제 전반의 신용 비용에 영향을 미칩니다.
        \item 은행은 지급준비금을 보유할지, 아니면 모기지와 같은 대출을 해줄지 결정합니다.
        \item 만약 연방기금 금리가 충분히 높다면, 은행은 새로운 모기지에 대해 높은 이자를 받아야 합니다. 그렇지 않으면 은행은 단순히 지급준비금을 보유하고 연방기금 시장에 빌려줄 수 있기 때문입니다.
        \item 따라서 연방기금 금리는 모기지 비용에 영향을 미칠 수 있습니다. 이 논리는 모기지뿐만 아니라 다른 모든 자산에도 적용됩니다. 모든 시장의 금리는 연방기금 금리와 연결되어 있습니다.
        \item 그러나 단기 만기 자산에 더 강한 영향을 미칠 것입니다. 단기 만기 자산은 연방기금 시장에 더 가까운 대체재이기 때문입니다. 이러한 대체 익일물 현금 시장은 연방기금 금리 변화에 즉시 영향을 받습니다.
    \end{itemize}
\end{enumerate}

\begin{infobox}[title=\textbf{중앙은행의 역할}]
미국 중앙은행인 연방준비제도(Federal Reserve)는 정책 목표를 달성하기 위해 연방기금 시장의 지급준비금 공급을 통제하여 연방기금 금리를 조절합니다. 이를 \textbf{공개 시장 조작 (Open Market Operations)}이라고 합니다.
\end{infobox}

\subsubsection*{핵심 요약: 연방기금 시장}
\begin{itemize}
    \item 은행은 법적으로 요구되는 지급준비금과 예상치 못한 유동성 수요를 충족시키기 위한 초과 지급준비금을 보유합니다.
    \item 연방기금 시장은 은행 시스템 전반에 걸쳐 현금을 재분배합니다.
    \item 연방기금 금리는 다른 모든 금리에 영향을 미칩니다.
\end{itemize}

\newpage
\subsection{환매조건부채권 (Repo) 시장 (The Repo Market)}
이 강의에서는 환매조건부채권(Repurchase Agreement) 시장을 논의합니다.

\subsubsection*{환매조건부채권 시장의 규모}
\begin{itemize}
    \item 환매조건부채권(Repo) 시장은 익일물 차입 및 대출의 핵심 시장입니다.
    \item \textbf{규모:} 2020년 말까지 4조 달러 이상의 환매조건부채권 대출이 미상환 상태였습니다.
    \item \textbf{인기 이유:} 환매조건부채권 시장은 가장 큰 현금 시장이며, 모든 유형의 금융 기관이 참여할 수 있습니다. 이는 은행만 접근할 수 있는 연방기금 시장과의 주요 차이점입니다.
    \item \textbf{주요 참여자:} 머니 마켓 뮤추얼 펀드(Money Market Mutual Funds)는 환매조건부채권 시장에서 가장 큰 대출 기관 중 하나입니다. 2019년 말, 이들은 약 1조 달러의 미상환 환매조건부채권을 보유하고 있었습니다.
\end{itemize}

\subsubsection*{환매조건부채권 시장의 작동 방식}
환매조건부채권 시장은 현금 시장입니다. 차입자는 대출자에게 익일물 현금을 빌립니다.

\begin{itemize}
    \item \textbf{담보 거래 (Secured Transaction):} 환매조건부채권 시장의 거래는 담보가 제공되는 거래입니다.
    \item \textbf{예시:} 오늘 청구서를 지불하기 위해 현금이 필요하고 국채(Treasury bond)를 가지고 있다고 가정해 봅시다. 내일은 현금이 풍부해질 것이므로 오늘 현금 부족을 겪고 있습니다.
    \begin{itemize}
        \item \textbf{대안 1 (비효율적):} 오늘 국채를 팔고, 청구서를 지불한 다음, 내일 새로운 국채를 살 수 있습니다. 그러나 채권을 사고파는 것은 비용이 많이 들 수 있습니다.
        \item \textbf{대안 2 (Repo):} 필요한 현금을 빌릴 수 있습니다. 환매조건부채권 거래에서는 국채를 현금과 교환하여 팔고, 내일 다시 사들이기로 약속합니다. 이를 \textbf{환매조건부채권 매도 (selling a repo)}라고 합니다.
    \end{itemize}
    \item \textbf{역할:}
    \begin{itemize}
        \item \textbf{환매조건부채권 매도자 (Repo Seller):} 다른 당사자로부터 현금을 빌리는 사람입니다.
        \item \textbf{환매조건부채권 매수자 (Repo Buyer):} 현금 대출을 제공하는 사람입니다. 매수자는 거래 기간 동안 국채의 새로운 소유자가 됩니다. 즉, 이는 담보 거래입니다.
        \item \textbf{이자:} 매수자는 보통 약간 더 높은 가격으로 국채를 다시 사들입니다. 이로써 매도자는 현금 대출에 대한 보상을 받습니다.
    \end{itemize}
\end{itemize}

\subsubsection*{환매조건부채권 거래의 위험과 헤어컷 (Haircut)}
\begin{itemize}
    \item \textbf{대출자의 주요 위험:} 차입자가 내일 국채를 다시 사들이지 못할 수 있습니다. 이 경우 대출자는 가치 없는 채권을 보유하게 될 수 있습니다. 이는 대출자가 담보를 보유하는 동안 담보 가치가 하락할 때 발생할 가능성이 더 높습니다.
    \item \textbf{구체적인 예시:}
    \begin{itemize}
        \item 당신이 오늘 200달러 가치의 국채를 팔고 내일 200.02달러에 다시 사들이기로 약속했다고 가정해 봅시다.
        \item 경제 상황이 변하여 시장 가격이 199달러로 떨어지면 어떻게 될까요? 당신은 약속대로 200.02달러에 다시 사들이거나, 약속을 어기고 공개 시장에서 199달러에 다른 국채를 살 수 있습니다.
        \item 대출자는 어떻게 될까요? 그들은 199달러 가치의 국채를 보유하고 있지만 200달러를 빌려주었습니다. 그들은 손실을 볼 가능성이 높습니다.
    \end{itemize}
    \item \textbf{헤어컷 (Haircut):} 이러한 위험을 줄이기 위해 환매조건부채권 거래에는 일반적으로 \textbf{헤어컷}이 적용됩니다. 대출자는 대출 금액에 헤어컷을 적용합니다. 담보의 전체 금액을 빌려주는 대신, 일부만 빌려줍니다.
    \item \textbf{헤어컷의 효과:} 헤어컷은 대부분의 환매조건부채권을 대출자에게 사실상 무위험하게 만듭니다.
    \item \textbf{예시 (헤어컷 적용):} 대출자가 2\%의 헤어컷을 부과한다고 가정해 봅시다. 따라서 차입자는 200달러 가치의 국채에 대해 196달러만 받습니다.
    \begin{itemize}
        \item 시장 가격이 199달러로 떨어져도 차입자는 여전히 채권을 다시 사들이고 싶어 할 것입니다 (199달러에 시장에서 사는 것보다 196달러에 다시 사는 것이 이득이므로).
        \item 채권 가격이 196달러 이상으로 유지된다면 대출자는 안전합니다.
    \end{itemize}
\end{itemize}

\subsubsection*{헤어컷의 크기를 결정하는 요인}
헤어컷의 크기는 담보와 관련된 세 가지 주요 요인에 의해 결정됩니다.

\begin{enumerate}
    \item \textbf{담보의 신용 품질 (Credit Quality):} 신용 위험이 낮을수록 헤어컷이 낮아집니다.
    \item \textbf{담보의 시장 가격 변동성 (Market Price Volatility):} 변동성이 낮을수록 헤어컷이 낮아집니다.
    \item \textbf{전반적인 시장 상황 (Overall Market Condition):} 담보를 쉽게 팔 수 있을수록 헤어컷이 낮아집니다.
\end{enumerate}

\subsubsection*{담보의 종류}
\begin{itemize}
    \item 다양한 자산이 환매조건부채권 거래에 사용될 수 있습니다. 그러나 가장 일반적인 담보는 \textbf{정부 보증 증권 (government backed securities)}입니다.
    \item \textbf{국채 (Treasuries)가 인기 있는 이유:}
    \begin{itemize}
        \item 신용 위험이 없습니다.
        \item 시장 가격이 안정적입니다.
        \item 거래하기 쉽습니다.
        \item 따라서 환매조건부채권 매수자에게 거의 위험을 주지 않으며 헤어컷도 낮습니다.
    \end{itemize}
    \item \textbf{일반 담보 환매조건부채권 (General Collateral Repo):} 정부 보증 증권을 포함하는 환매조건부채권 거래는 위험이 낮습니다. 이 경우 대출자는 특정 담보 유형에 대해 덜 까다롭습니다. 담보가 지정되기 전에 합의되는 환매조건부채권 거래를 일반 담보 환매조건부채권이라고 합니다. 여기서 일반 담보는 적격 정부 보증 증권 목록을 의미합니다. 일반 담보 환매조건부채권은 헤어컷이 가장 낮으며 가장 일반적인 유형의 환매조건부채권 거래입니다.
\end{itemize}

\subsubsection*{핵심 요약: 환매조건부채권 시장}
\begin{itemize}
    \item 환매조건부채권은 단기 담보 대출의 한 형태입니다.
    \item 환매조건부채권 시장은 은행뿐만 아니라 모든 금융 기관에게 중요한 현금 시장입니다.
    \item 헤어컷은 대출자의 손실 위험을 줄입니다.
\end{itemize}

\newpage
\subsection{단기 금리 비교 (Short-term Interest Rates)}
이 강의에서는 현금 시장에서 중요한 두 가지 단기 금리를 소개합니다. 이 금리들은 모기지 및 신용카드 대출 비용을 통해 가계에, 사업 대출 비용을 통해 기업에 영향을 미칩니다.

\subsubsection*{런던 은행 간 대출 금리 (LIBOR)}
\begin{itemize}
    \item \textbf{정의:} LIBOR는 런던 은행 간 시장의 금리입니다. 이 시장은 은행들이 서로 돈을 빌려주고 빌려주는 시장입니다.
    \item \textbf{통화:} 대부분의 대출은 미국 달러로 이루어지지만, 은행들은 다른 통화도 빌릴 수 있습니다. 유럽에서 이루어지는 미국 달러 대출은 유로달러 대출(Euro dollar loans)이라고도 불립니다.
    \item \textbf{특징:}
    \begin{itemize}
        \item 이 시장의 대출은 \textbf{무담보 (unsecured)}입니다. 즉, 담보가 필요하지 않습니다.
        \item 담보가 필요 없기 때문에, 이 은행 간 대출은 \textbf{신용 위험 (credit risk)}에 노출됩니다. 은행이 대출금을 상환하지 못할 가능성이 있습니다.
        \item \textbf{만기:} 익일물(overnight), 1주, 1개월, 2개월, 3개월, 6개월, 12개월의 7가지 만기가 존재합니다.
    \end{itemize}
    \item \textbf{계산 방식:} LIBOR는 실제 거래보다는 \textbf{설문 조사 (surveys)}를 기반으로 계산됩니다. 매일 가장 크고 거래량이 많은 은행들을 대상으로 설문 조사를 실시하며, LIBOR는 이 설문 조사 결과에 기반한 평균 금리입니다.
\end{itemize}

\begin{examplebox}[title=\textbf{LIBOR와 연방기금 금리의 관계}]
익일물 미국 달러 LIBOR는 최근 몇 년간 미국 은행 간 금리인 유효 연방기금 금리(Effective Federal Funds Rate)와 유사한 패턴을 보입니다. 이는 미국과 해외에서 미국 달러를 빌리는 비용이 매우 밀접하게 움직인다는 것을 의미합니다.

\textbf{왜 그럴까요?} 만약 LIBOR가 더 낮다면, 은행들은 런던에서 달러를 빌려 미국에 빌려줌으로써 이익을 얻을 수 있습니다. 이러한 기본적인 \textbf{차익 거래 (arbitrage)}는 LIBOR와 연방기금 금리가 대략적으로 같도록 보장합니다.
\end{examplebox}

\subsubsection*{LIBOR와 광범위한 경제의 연결}
\begin{itemize}
    \item LIBOR는 은행의 자금 조달 비용을 측정하므로, 기업 대출과 같은 많은 다른 대출의 \textbf{기준 금리 (reference rate)}가 되었습니다.
    \item 이러한 대출은 종종 \textbf{변동 금리 (floating rate)}입니다. 이는 이자율이 변동 기준 금리(예: LIBOR)에 고정된 가산금리(mark up)를 더한 것과 같다는 의미입니다.
    \item \textbf{예시:} 월마트의 2019년 규제 서류에 따르면, 2020년 5월부터 2024년 5월 사이에 만기가 도래하는 약정 신용 한도는 일반적으로 LIBOR + 10bp에서 LIBOR + 75bp 사이의 이자율을 적용받습니다.
    \item 2020년까지 LIBOR는 전 세계적으로 250조 달러 이상의 대출을 뒷받침했습니다.
\end{itemize}

\subsubsection*{LIBOR의 문제점: 조작 논란과 대체}
\begin{itemize}
    \item 2000년부터 2020년까지의 기간을 보면, LIBOR가 여러 차례 급등하여 연방기금 금리를 따르지 않는 경우가 있었습니다.
    \item \textbf{원인 1: 신용 위험:} LIBOR는 무담보 대출이므로, 런던 시장 은행들의 건전성에 의문이 제기되면 은행 간 금리 사이에 격차가 발생할 수 있습니다.
    \item \textbf{원인 2: 조작 논란:} 2008년 금융 위기 이후, LIBOR 조작 의혹이 제기되었습니다. LIBOR는 가장 큰 은행들의 금리 설문 조사를 기반으로 하며, 이 설문 조사는 정직성을 요구합니다. 그러나 이 설문 조사가 실제로 조작되었다는 증거가 나왔습니다. 총 90억 달러 이상의 벌금이 은행들에 부과되었습니다.
    \item \textbf{대체:} 추가 조작을 방지하기 위해 규제 당국은 LIBOR를 대출의 기준 금리에서 제외하도록 요구했습니다. 2022년부터 새로운 금리가 의무화되었습니다.
\end{itemize}

\subsubsection*{담보부 익일물 자금 조달 금리 (SOFR)}
\begin{itemize}
    \item \textbf{정의:} 새로운 금리는 \textbf{담보부 익일물 자금 조달 금리 (Secured Overnight Financing Rate, SOFR)}라고 불립니다. 뉴욕 연방준비은행이 매일 발표합니다.
    \item \textbf{특징:}
    \begin{itemize}
        \item SOFR 금리는 국채를 담보로 하는 익일물 환매조건부채권(repo) 거래의 금리입니다.
        \item SOFR은 설문 조사가 아닌 \textbf{실제 거래 (transactions)}를 기반으로 합니다.
        \item SOFR은 \textbf{담보부 금리 (secured rate)}이므로 \textbf{무위험 (risk free)}으로 간주됩니다.
        \item 반면 LIBOR는 무담보 금리이므로 일부 신용 위험을 포함합니다.
    \end{itemize}
    \item \textbf{규모:} 2020년 기준으로 SOFR은 약 1조 달러 규모의 일일 거래 시장을 기반으로 합니다. LIBOR는 이 규모의 절반 미만에 기반했습니다.
\end{itemize}

\subsubsection*{핵심 요약: 단기 금리}
\begin{itemize}
    \item LIBOR와 SOFR은 단기 금리이며, 각각 무담보 및 담보 대출에 대한 은행의 자금 조달 비용을 나타냅니다.
    \item LIBOR는 신용 위험에 노출되어 있으며 과거에 조작된 사례가 있습니다.
    \item 규제 당국은 무위험 금리인 SOFR이 LIBOR를 기준 금리로 대체하도록 의무화했습니다.
\end{itemize}

\newpage
\section{장기 금리 (Long-term Rates)}

\subsection{명목 금리와 실질 금리 (Nominal and Real Interest Rates)}
이 강의에서는 명목 금리와 실질 금리를 논의하고, 이 두 금리를 인플레이션율과 연결하는 피셔 방정식(Fisher Equation)을 살펴봅니다. 마지막으로, 실질 금리를 지급하는 몇 안 되는 증권 중 하나인 물가연동국채(Treasury Inflation Protected Securities, TIPS)를 사례 연구로 살펴봅니다.

\subsubsection*{명목 수익률과 실질 수익률의 차이}
투자에 대한 명목 수익률과 실질 수익률의 차이를 예시를 통해 이해할 수 있습니다.

\begin{itemize}
    \item \textbf{예시:} 은행 계좌에 100달러가 있고, 연 10\%의 이자를 지급한다고 가정해 봅시다.
    \begin{itemize}
        \item 1년 후, 10달러를 얻어 총 110달러가 됩니다.
        \item 은행이 지급하는 이 10달러가 100달러 투자에 대한 \textbf{명목 수익률 (nominal return)}입니다. 이 10\%의 이자율을 \textbf{명목 금리 (nominal interest rate)}라고 합니다.
    \end{itemize}
    \item \textbf{실질 수익률 (real return):} 이 이익이 1년 동안 구매하고 소비할 수 있는 실제 상품의 가치로 얼마나 되는가? 이것이 투자에 대한 실질 수익률이 의미하는 바입니다.
\end{itemize}

\begin{examplebox}[title=\textbf{명목 금리와 실질 금리 예시: 주유비}]
\begin{itemize}
    \item \textbf{현재:} 오늘 휘발유 1갤런을 4달러에 살 수 있습니다. 따라서 100달러는 25갤런의 휘발유 가치입니다.
    \item \textbf{인플레이션 가정:} 인플레이션, 즉 경제 전반의 물가 상승률이 5\%라고 가정해 봅시다. 내년에는 휘발유 가격이 4.20달러로 오릅니다. 같은 25갤런의 휘발유를 구매하려면 105달러가 필요합니다.
    \item \textbf{투자 수익률:}
    \begin{itemize}
        \item \textbf{명목 수익률:} 은행은 여전히 저축에 대해 10달러를 지급하므로, 명목 수익률은 10달러입니다.
        \item \textbf{실질 수익률:} 휘발유 가격이 올랐습니다. 25갤런의 휘발유를 구매한 후에는 5달러가 남습니다. 따라서 인플레이션을 고려한 실질 수익률은 5달러입니다.
        \item 명목 수익률 10달러에서 물가 상승분 5달러를 빼면 5달러의 실질 수익률이 나옵니다.
    \end{itemize}
\end{itemize}
이 예시는 명목 금리와 실질 금리가 인플레이션에 의해 연결되어 있음을 보여줍니다.
\end{examplebox}

\subsubsection*{피셔 방정식 (Fisher Equation)}
일반적으로 \textbf{실질 금리 (real interest rate)}는 \textbf{명목 금리 (nominal interest rate)}에서 \textbf{인플레이션 (inflation)}을 뺀 값과 거의 같습니다.

\begin{infobox}[title=\textbf{피셔 방정식}]
명목 금리 $i$는 실질 금리 $r$에 인플레이션 $\pi$를 더한 값과 대략적으로 같습니다.
$$ i \approx r + \pi $$
이 방정식은 미국의 경제학자 어빙 피셔(Irving Fisher)의 이름을 따서 명명되었습니다.
\end{infobox}

\subsubsection*{사례 연구: 물가연동국채 (TIPS)}
물가연동국채(Treasury Inflation Protected Securities, TIPS)는 인플레이션을 명시적으로 고려하는 채권 투자입니다.

\begin{itemize}
    \item \textbf{작동 방식:} TIPS는 채권의 원금 가치가 인플레이션율에 따라 증가합니다. 채권에 지급되는 이자는 원금 가치의 백분율이므로, TIPS에 지급되는 이자도 인플레이션에 따라 증가합니다. 따라서 TIPS의 이자율은 \textbf{실질 금리}입니다.
    \item \textbf{구체적인 예시:}
    \begin{itemize}
        \item 100달러짜리 TIPS 채권을 구매하고 연 5\%의 이자를 지급한다고 가정해 봅시다.
        \item 올해 인플레이션이 2\%라고 가정하면, 내년에 TIPS 채권의 원금은 2\% 증가하여 102달러가 됩니다.
        \item \textbf{채권의 연간 총 수익률:} 5\% 이자 + 2\% 원금 가치 상승 = 총 7\%의 명목 수익률.
        \item \textbf{실질 수익률:} 피셔 방정식에 따르면, 실질 수익률은 명목 수익률에서 인플레이션을 뺀 값입니다. 2\%의 원금 가치 상승은 2\%의 인플레이션을 상쇄합니다. 따라서 실질 수익률은 5\%입니다.
        \item TIPS는 일정한 실질 수익률을 제공하도록 설계되었습니다.
    \end{itemize}
\end{itemize}

\subsubsection*{TIPS 시장의 중요성: 미래 인플레이션 기대}
경제학자들은 TIPS 시장에 왜 주목할까요?

\begin{itemize}
    \item 오늘 투자를 할 때, 우리는 미래 인플레이션이 어떻게 될지 알지 못합니다. 인플레이션은 실질 수익률에 중요합니다.
    \item TIPS 채권의 가격을 사용하여 시장의 \textbf{미래 인플레이션 기대 (market's expectations of future inflation)}를 파악할 수 있습니다.
    \item \textbf{작동 방식:}
    \begin{itemize}
        \item TIPS는 일정한 실질 수익률을 제공하며 인플레이션 보호 기능을 제공합니다.
        \item 예상 인플레이션이 높으면 인플레이션 보호의 가치가 커집니다. TIPS 채권은 인플레이션 보호를 제공하므로, TIPS의 현재 시장 가격은 상승할 것입니다. 가격이 높으면 TIPS의 수익률은 낮아집니다.
        \item 반대로, 인플레이션이 낮을 것으로 예상되면 인플레이션 보호의 가치가 떨어집니다. 이 경우 TIPS 채권 가격은 하락하고, TIPS 채권의 수익률은 상승합니다.
        \item 따라서 TIPS 가격 또는 수익률은 인플레이션 기대를 포착합니다. 실제로 TIPS 수익률과 일반 국채 수익률을 비교하여 예상되는 미래 인플레이션을 측정할 수 있습니다.
    \end{itemize}
\end{itemize}

\begin{examplebox}[title=\textbf{TIPS 수익률 데이터 분석}]
5년 만기 TIPS 채권 수익률과 5년 만기 국채 수익률 데이터를 보면, 두 수익률 사이에 차이가 나는 경우가 있습니다. 수익률 차이가 크다는 것은 투자자들이 인플레이션 보호에 기꺼이 비용을 지불한다는 의미이며, 이는 투자자들이 5년 후 높은 미래 인플레이션을 예상한다는 것을 뜻합니다.

\begin{itemize}
    \item \textbf{2009년:} TIPS 수익률이 높았습니다. 투자자들은 이 시기에 더 높은 인플레이션을 예상했습니다.
    \item \textbf{2013년 초 및 2021년:} 반대의 상황이었습니다. 투자자들은 낮은 미래 인플레이션을 예상했습니다.
\end{itemize}
물론 시장의 예측이 결국 틀릴 수도 있지만, 이는 TIPS 가격에 반영된 기대치였습니다.
\end{examplebox}

\subsubsection*{핵심 요약: 명목 금리와 실질 금리}
\begin{itemize}
    \item 명목 금리는 상품 가격 변화를 고려하지 않습니다.
    \item 명목 금리는 실질 금리 + 인플레이션과 같습니다. 이것이 피셔 방정식입니다.
    \item TIPS 채권은 인플레이션 기대를 포착합니다.
\end{itemize}

\newpage
\subsection{국채와 무위험 금리 (Treasury Securities and The Risk-Free Rate)}
이 강의에서는 대표적인 무위험 자산인 미국 국채(US Treasury Securities)를 논의합니다.

\subsubsection*{미국 국채 시장의 특징}
\begin{itemize}
    \item \textbf{세계 최대 채권 시장:} 2020년 기준으로 미국 국채 시장은 전 세계에서 가장 큰 채권 시장입니다. 당시 미상환 증권 가치는 24조 달러를 넘었습니다.
    \item \textbf{다양한 투자자:} 거의 모든 유형의 미국 금융 기관(은행, 뮤추얼 펀드, 보험 회사, 연기금, 연방준비제도 등)이 국채를 소유하고 있습니다. 또한 국채의 약 3분의 1은 국제 투자자들이 소유하고 있습니다.
\end{itemize}

\subsubsection*{미국 국채의 종류}
미국 국채는 만기에 따라 세 가지 범주로 나뉩니다. 이들은 평생 고정 이자 지급을 하는 정부 채권입니다.

\begin{enumerate}
    \item \textbf{재무부 단기채 (Treasury Bills, T-Bills):}
    \begin{itemize}
        \item 만기가 가장 짧습니다 (4주에서 52주).
        \item 미국 국채 증권의 약 20\%를 차지합니다.
    \end{itemize}
    \item \textbf{재무부 중기채 (Treasury Notes, T-Notes):}
    \begin{itemize}
        \item 중간 만기 범위의 증권입니다 (2년에서 10년).
        \item 미상환 미국 국채 증권의 대부분을 차지합니다.
    \end{itemize}
    \item \textbf{재무부 장기채 (Treasury Bonds, T-Bonds):}
    \begin{itemize}
        \item 만기가 20년에서 30년인 점을 제외하고는 재무부 중기채와 본질적으로 동일합니다.
        \item '장기 채권 (long bond)'이라고도 불립니다.
    \end{itemize}
\end{enumerate}

\subsubsection*{미국 국채가 투자자들에게 인기 있는 이유}
\begin{itemize}
    \item \textbf{신용 위험 없음 (Credit Risk-Free):} 미국 국채는 신용 위험이 없는 것으로 간주됩니다. 즉, 투자자들은 만기까지 국채를 보유하는 한 이자와 원금 지급을 보장받습니다. 이러한 지급은 미국 정부의 완전한 신뢰와 신용(full faith and credit)으로 보장됩니다.
    \item \textbf{담보 활용:} 미국 국채는 신용 위험이 없기 때문에 차입 약정에서 담보로 자주 사용됩니다. 예를 들어, 수조 달러 규모의 환매조건부채권(repo) 시장에서 담보로 활용됩니다.
    \item \textbf{무위험 금리 (Risk-Free Interest Rate)의 척도:} 국채의 이자율은 무위험 금리의 좋은 척도로 간주됩니다. 무위험 금리는 전 세계 많은 금융 계약에서 벤치마크 이자율로 사용됩니다.
\end{itemize}

\begin{infobox}[title=\textbf{다른 선진국 국채보다 미국 국채가 선호되는 이유}]
독일, 중국, 일본과 같이 채무 불이행 위험이 없는 다른 선진국의 국채보다 미국 국채가 선호되는 세 가지 주요 이유가 있습니다.
\begin{enumerate}
    \item \textbf{시장 규모와 유동성:} 미국 국채 시장은 다른 어떤 시장보다 훨씬 큽니다. 언제든지 다양한 투자자들이 국채 구매에 관심을 가지므로, 거래 비용이 거의 없이 쉽게 매도할 수 있습니다. 시장에서 증권이 적은 비용으로 쉽게 거래될 수 있을 때, 그 시장을 \textbf{유동적 (liquid)}이라고 합니다. 미국 국채 시장은 세계에서 가장 유동적인 증권 시장입니다.
    \item \textbf{달러 자산 보유 수단:} 전 세계 대부분의 금융 거래는 미국 달러로 이루어집니다. 미국 국채를 보유하는 것은 미국 달러를 보유하는 한 가지 방법입니다. 이러한 이유로 미국 국채는 종종 \textbf{현금성 자산 (cash-like assets)}이라고 불립니다. 2020년 말까지 기업과 가계는 미국 국채에만 투자하는 머니 마켓 뮤추얼 펀드에 2조 달러 이상의 현금을 예치했습니다.
    \item \textbf{중앙은행의 달러 준비금:} 전 세계 중앙은행들이 달러 준비금을 보유하는 가장 일반적인 방법이 미국 국채입니다.
\end{enumerate}
\end{infobox}

\subsubsection*{미국 국채 투자에 항상 위험이 없을까?}
만기까지 국채를 보유하면 지급을 받으므로 이 경우에는 위험이 없습니다. 그러나 만기 전에 매도해야 하는 경우 손실을 볼 수 있습니다. 이는 채권의 현재 시장 가치가 하락할 때 발생합니다. 국채 투자자들은 두 가지 주요 위험 원인에 대해 우려합니다.

\begin{enumerate}
    \item \textbf{인플레이션 (Inflation):}
    \begin{itemize}
        \item \textbf{예시:} 액면가 100달러, 명목 이자율 2\%인 52주 만기 T-Bill을 구매했다고 가정해 봅시다. 인플레이션이 2\%일 것으로 예상했다고 가정합니다. 만기 시 100달러의 액면가와 2달러의 이자를 받습니다. 동시에 T-Bill을 구매했을 때 100달러였던 상품은 인플레이션으로 인해 102달러가 됩니다. 벌어들인 이자로 T-Bill을 구매했을 때 살 수 있었던 것과 동일한 상품을 여전히 살 수 있습니다.
        \item \textbf{예상보다 높은 인플레이션:} 올해 인플레이션이 5\%라고 가정해 봅시다. 만기 시 여전히 100달러의 액면가와 2달러의 이자를 받습니다. 그러나 이제 상품은 더 높은 인플레이션으로 인해 105달러가 됩니다. 이는 이전과 동일한 상품을 살 수 없다는 의미입니다.
        \item \textbf{결과:} 명목 수익률은 2\%였지만, 인플레이션을 조정한 실질 수익률은 마이너스 3\%가 됩니다. 인플레이션이 예상보다 높으면 국채의 시장 가치는 하락합니다. 현재 관점에서 국채는 구매력이 적으므로 가치가 떨어집니다.
    \end{itemize}
    \item \textbf{금리 변화 (Changes in the Interest Rate):}
    \begin{itemize}
        \item \textbf{예시:} 100달러에 2\% 이자를 지급하는 2년 만기 국채를 구매했다고 가정해 봅시다. 1년 후 이 투자를 팔고 싶다고 가정합니다. 동시에 이 투자에는 1년 만기가 남아 있습니다. 이 채권으로 얼마를 받을 수 있을까요?
        \item \textbf{결과:} 이는 가장 가까운 대체재인 새로 발행된 52주 만기 T-Bill의 이자율에 따라 달라집니다. 새로 발행된 52주 만기 T-Bill의 이자율이 더 높다면(예: 3\%), 이는 2\%만 지급하는 당신의 채권보다 더 매력적인 투자입니다. 새로운 T-Bill이 100달러라면, 당신의 투자는 100달러 미만이 될 것입니다.
        \item \textbf{결론:} 금리가 상승하면 이전에 발행된 장기 국채의 시장 가치는 하락할 수 있습니다. 장기 국채는 채권의 수명 동안 금리가 변동할 위험이 더 큽니다. 이 위험에 대한 보상을 \textbf{기간 프리미엄 (term premium)}이라고 합니다.
    \end{itemize}
\end{enumerate}

\subsubsection*{장기 국채 명목 금리의 구성 요소}
이러한 위험은 장기 국채의 명목 금리에 반영됩니다. 세 가지 구성 요소가 있습니다.

\begin{enumerate}
    \item \textbf{실질 금리 (Real Interest Rate):} 인플레이션을 고려한 금리입니다.
    \item \textbf{예상 인플레이션율 (Expected Rate of Inflation):} 미래에 예상되는 물가 상승률입니다.
    \item \textbf{기간 프리미엄 (Term Premium):} 금리 변화를 포함한 장기 위험을 포착합니다.
\end{enumerate}

\subsubsection*{핵심 요약: 국채와 무위험 금리}
\begin{itemize}
    \item 미국 국채 시장은 세계에서 가장 크고 유동적인 채권 시장입니다.
    \item 미국 국채는 신용 위험이 없습니다. 이는 그 이자율이 무위험 금리의 좋은 대용치라는 것을 의미합니다.
    \item 미국 국채의 장기 명목 금리는 실질 금리, 예상 인플레이션, 기간 프리미엄으로 분해될 수 있습니다.
\end{itemize}

\newpage
\subsection{금리 기간 구조 (Term Structure of Interest Rates)}
이 강의에서는 금리 기간 구조가 무엇인지 논의합니다. 금리 기간 구조가 거시 경제에 대한 중요한 정보를 제공하는 이유를 살펴봅니다. 특히, 경기 침체 전에 금리 기간 구조가 어떻게 행동하는지 알아봅니다.

\subsubsection*{금리 기간 구조의 정의}
\begin{itemize}
    \item \textbf{정의:} \textbf{금리 기간 구조 (Term Structure of Interest Rates)}는 특정 시점에 동일한 발행 주체가 발행한 만기가 다른 증권의 금리를 비교하는 것을 의미합니다.
    \item \textbf{수익률 (Yields) 사용:} 금리 기간 구조에는 채권 소유자에게 지급되는 쿠폰 이자율(coupon interest rates)이 아니라, 채권 가격을 고려한 \textbf{수익률 (yields)}을 사용합니다.
    \item \textbf{예시:} 액면가 100달러에 5\% 이자를 지급하는 채권이 있다고 가정해 봅시다.
    \begin{itemize}
        \item 채권을 99달러에 구매하면 여전히 5달러의 이자를 받습니다. 현재 수익률은 5/99, 즉 5.051\%입니다.
        \item 반대로, 같은 5\% 이자 지급 채권을 101달러에 구매하면 수익률은 5/101, 즉 4.95\%가 됩니다.
    \end{itemize}
    \item \textbf{채권 가격과 수익률의 역관계:} 채권 수익률과 채권 가격은 반대 방향으로 움직입니다. 채권 가격이 상승하면 수익률은 하락하고, 채권 가격이 하락하면 수익률은 상승합니다.
\end{itemize}

\subsubsection*{수익률 곡선 (Yield Curve)}
\begin{itemize}
    \item 우리는 매일 수익률의 변화와 그에 따른 금리 기간 구조를 관찰할 수 있으며, 이를 일반적으로 \textbf{수익률 곡선 (Yield Curve)}이라고 합니다.
    \item \textbf{가장 중요한 수익률 곡선:} 미국 국채(US Treasury Securities)를 기반으로 합니다.
    \item \textbf{미국 국채가 좋은 척도가 되는 세 가지 이유:}
    \begin{enumerate}
        \item 시간이 지남에 따라 변할 수 있는 신용 위험이 없습니다.
        \item 다양한 만기로 발행됩니다.
        \item 증권 가격과 수익률을 왜곡할 수 있는 거래 비용이 거의 없습니다.
    \end{enumerate}
\end{itemize}

\begin{examplebox}[title=\textbf{실제 수익률 곡선 예시: 2021년 6월 1일}]
2021년 6월 1일의 미국 국채 수익률 곡선을 보면, 단기 만기 채권의 수익률은 거의 0\%에 가깝습니다. 1년 만기 미국 국채도 연 0.04\%에 불과합니다. 그러나 10년 만기 국채는 연 1.62\%의 수익률을 보입니다.

\textbf{왜 단기 만기와 장기 만기 수익률에 차이가 있을까요?}
원금을 갚는 데 10년이 걸리는 프로젝트에 자금을 조달하는 두 가지 대안을 생각해 봅시다.
\begin{enumerate}
    \item \textbf{대안 1:} 10년 만기 채권을 발행합니다.
    \item \textbf{대안 2:} 1년 만기 채권을 발행하고, 1년 후 새로운 1년 만기 채권으로 부채를 롤오버(roll over)합니다. 이 경우 9번 더 부채를 롤오버해야 합니다.
\end{enumerate}
만약 1년 만기 채권으로 부채를 롤오버할 때 항상 처음과 동일한 수익률을 얻을 수 있다면, 10번의 1년 만기 채권 발행 비용은 10년 만기 채권 발행 비용과 같을 것이고, 따라서 수익률도 같을 것입니다. 만약 같지 않다면, 더 저렴한 대안을 선택하는 것이 유리할 것입니다.

2021년 6월 1일 수익률 곡선에서는 단기 채권의 수익률이 더 낮습니다. 이는 시장 참여자들이 향후 10년 동안 어느 시점에서 금리가 상승할 것으로 예상한다는 의미입니다. 즉, 1년 만기 채권을 롤오버할 때 처음 발행했던 1년 만기 채권보다 더 높은 수익률을 제공해야 할 것으로 예상하는 것입니다.
\end{examplebox}

\subsubsection*{높은 금리 기대를 유발하는 요인}
\begin{enumerate}
    \item \textbf{연방준비제도의 단기 금리 인상 기대:} 시장 참여자들은 연방준비제도가 단기 금리를 인상할 것으로 예상할 수 있습니다. 이는 일반적으로 경제 확장기(economic expansions)에 예상되며, 이때 수익률 곡선은 일반적으로 \textbf{우상향 (upward sloping)}합니다.
    \item \textbf{높은 인플레이션 기대:} 시장 참여자들은 더 높은 인플레이션을 예상할 수 있습니다. 미국 국채의 금리는 명목 금리이므로, 투자자들을 계속 유치하기 위해서는 장기 만기 국채의 수익률이 인플레이션 증가를 상쇄하기 위해 증가해야 합니다.
    \item \textbf{재정 정책으로 인한 정부 부채 증가:} 정부 부채를 증가시키는 재정 정책은 미국 국채의 공급을 늘릴 것입니다. 이 경우, 투자자들이 추가 공급되는 미국 국채를 구매하도록 유인하기 위해 장기 만기 국채의 수익률이 증가해야 합니다.
\end{enumerate}

\subsubsection*{수익률 곡선과 경기 침체 예측}
\begin{itemize}
    \item \textbf{하향 경사 (Downward Sloping) 또는 역전된 (Inverted) 수익률 곡선:} 2019년 9월 4일의 수익률 곡선을 보면, 곡선이 앞쪽에서 하향 경사 또는 역전되어 있습니다. 이는 시장 참여자들이 미래에 금리가 하락할 것으로 예상해야 한다는 의미입니다.
    \item \textbf{언제 금리 하락을 예상할까요?} 바로 \textbf{경기 침체 (recession)} 전입니다.
    \item \textbf{10년물과 2년물 국채 수익률 차이:} 지난 40년간 10년 만기 국채와 2년 만기 국채 수익률의 차이를 보면, 거의 항상 이 차이가 양수였습니다. 즉, 10년물 국채 수익률이 2년물 국채 수익률보다 높았습니다. 대부분의 시간 동안 시장 참여자들은 금리가 상승할 것으로 예상했으며, 이는 경제 확장과 일치합니다.
    \item \textbf{경기 침체 전의 패턴:} 그림의 회색 막대는 경기 침체를 나타냅니다. 경기 침체기에는 수요가 낮으므로 금리가 하락하는 경향이 있습니다. 회색 막대 바로 직전에 수익률 곡선이 어떻게 되는지 주목하세요. 1980년대, 1990년대 초, 2000년 닷컴 버블, 2008년 금융 위기, 그리고 코로나19 팬데믹 이전의 경기 침체 직전에 매번 역전되었습니다.
\end{itemize}

\begin{summarybox}[title=\textbf{수익률 곡선의 중요성}]
요약하자면, 국채 수익률 곡선은 경제의 후속 성과를 강력하게 예측하는 지표이며, 따라서 시장 참여자와 정책 입안자들이 면밀히 주시합니다.
\end{summarybox}

\subsubsection*{핵심 요약: 금리 기간 구조}
\begin{itemize}
    \item 금리 기간 구조 또는 수익률 곡선은 만기가 다른 채권들의 수익률을 비교합니다.
    \item 수익률 곡선은 면밀히 주시되는 경제 지표입니다.
    \item \textbf{우상향 수익률 곡선:} 시장 참여자들이 미래에 금리가 상승할 것으로 예상하며, 이는 경제 확장기에 발생합니다.
    \item \textbf{역전된 하향 경사 수익률 곡선:} 미래 금리 하락을 시사하며, 일반적으로 경제 확장의 끝을 알리는 신호입니다.
\end{itemize}

\newpage
\section{기타 자금 시장 (Other Funding Markets)}

\subsection{머니 마켓 뮤추얼 펀드 (Money Market Mutual Funds)}
이 강의에서는 머니 마켓 뮤추얼 펀드(Money Market Mutual Funds)를 자세히 논의합니다. 이 펀드들은 경제에서 특히 은행과 비은행 기관에 단기 자금을 제공하는 데 중요한 역할을 합니다.

\subsubsection*{머니 마켓 펀드의 규모와 인기}
\begin{itemize}
    \item \textbf{규모:} 2021년 기준으로 머니 마켓 펀드는 4조 달러 이상의 자산을 운용했습니다. 이는 기업과 가계의 총 은행 예금보다 많은 금액입니다.
    \item \textbf{코로나19 팬데믹의 영향:} 그래프를 보면 코로나19 팬데믹의 영향으로 저축이 급증한 것을 명확히 알 수 있습니다.
    \item \textbf{인기 이유:} 머니 마켓 펀드는 기업과 가계를 위한 저축 수단입니다.
    \item \textbf{탄생 배경:} 1970년대 미국에서 은행 예금에 지급할 수 있는 이자율을 제한하는 은행 규제인 \textbf{규제 Q (Regulation Q)} 때문에 탄생했습니다.
    \begin{itemize}
        \item 규제 Q에 따라 은행이 고객에게 지급할 수 있는 최대 이자율은 5.25\%로 제한되었습니다.
        \item 1970년대와 1980년대에는 인플레이션과 시장 금리가 5\%보다 훨씬 높았기 때문에 은행 저축 계좌는 매력이 없었습니다.
        \item 머니 마켓 펀드는 이 규제를 받지 않아 더 높은 이자율을 지급할 수 있었습니다.
    \end{itemize}
\end{itemize}

\subsubsection*{머니 마켓 펀드의 사업 모델}
\begin{itemize}
    \item \textbf{투자 프로필:} 머니 마켓 펀드는 초안전 단기 채무 증권에 투자합니다. 이러한 투자는 신용 위험이 거의 없습니다.
    \item \textbf{주요 투자처 (2021년 기준):}
    \begin{itemize}
        \item 미국 국채(US Treasury Securities)
        \item 미국 국채를 담보로 하는 환매조건부채권(Repurchase Agreements)
    \end{itemize}
\end{itemize}

\subsubsection*{머니 마켓 펀드의 종류}
모든 머니 마켓 펀드를 포함하는 이전 차트와 달리, 일부 펀드는 특정 투자에 특화되어 있습니다.

\begin{itemize}
    \item \textbf{국채 펀드 (Treasury Funds):} 오직 미국 국채에만 투자하는 머니 마켓 펀드입니다.
    \item \textbf{정부 펀드 (Government Funds):} 다른 정부 증권에도 투자하는 펀드입니다.
    \item \textbf{특징:} 이들은 가장 안전한 투자로 간주되지만, 수익률도 가장 낮습니다.
    \item \textbf{프라임 펀드 (Prime Funds):} 다른 단기 채무 증권에도 투자하는 머니 마켓 펀드입니다.
    \begin{itemize}
        \item 기업이 발행하는 기업어음(commercial paper)
        \item 은행이 발행하는 양도성 예금증서(certificates of deposit)
        \item 헤지 펀드가 발행하는 환매조건부채권(purchase agreements)
    \end{itemize}
    \item \textbf{특징:} 이러한 투자는 약간 더 위험하지만, 수익률도 약간 더 높습니다.
\end{itemize}

\subsubsection*{머니 마켓 펀드와 은행 예금의 비교}
머니 마켓 펀드는 뮤추얼 펀드의 한 종류입니다. 현금을 머니 마켓 계좌에 예치하면 뮤추얼 펀드의 주식을 구매하는 것입니다. 다른 투자 펀드와 달리 머니 마켓 펀드는 초안전 투자에 집중하고 손실을 피하려고 노력합니다.

\begin{itemize}
    \item 투자자는 오늘 1달러에 머니 마켓 펀드 주식을 구매하면 내일 1달러와 이자를 받고 상환할 수 있을 것으로 기대합니다.
    \item 머니 마켓 펀드 주식은 요구불 상환(redeemable on demand)이 가능합니다. 언제든지 돈을 인출할 수 있습니다.
    \item 이러한 특징 때문에 머니 마켓 펀드 투자는 은행 예금과 매우 유사하게 보입니다.
\end{itemize}

그러나 머니 마켓 펀드와 은행 예금 사이에는 두 가지 중요한 차이점이 있습니다.

\begin{enumerate}
    \item \textbf{이자율:} 머니 마켓 펀드는 거의 항상 더 높은 이자율을 지급합니다. 대규모 급여를 관리하는 기업의 경우, 작은 이자율 차이도 큰 금액을 의미할 수 있습니다.
    \item \textbf{예금 보험 (Deposit Insurance):}
    \begin{itemize}
        \item 은행이 파산하면 정부 예금 보험은 예금자들이 모든 돈을 돌려받을 수 있도록 보장합니다.
        \item 머니 마켓 펀드 투자는 정부에 의해 보험에 가입되어 있지 않습니다. 특히 프라임 펀드의 경우 작은 손실 위험이 있습니다.
    \end{itemize}
\end{enumerate}

\subsubsection*{머니 마켓 펀드의 중요성과 문제 발생 시 영향}
머니 마켓 펀드는 금융 시스템의 핵심입니다. 이들은 단기 자금 시장의 초석이며, 광범위한 금융 기관과 기업에 대규모 현금을 대출해 줍니다. 이러한 자금 조달은 단기간에 대체하기 어렵습니다.

\begin{itemize}
    \item 머니 마켓 펀드가 문제가 생기는 경우는 드물지만, 금융 시장과 광범위한 경제에 매우 파괴적일 수 있습니다.
\end{itemize}

\subsubsection*{사례 연구: 2008년 금융 위기 중 리저브 프라이머리 펀드 (Reserve Primary Fund)}
\begin{itemize}
    \item 2008년 기준으로 리저브 프라이머리 펀드는 미국에서 가장 오래되고 큰 머니 마켓 펀드 중 하나였습니다.
    \item 이 펀드는 투자 은행인 리먼 브라더스(Lehman Brothers)의 기업어음에 크게 투자했습니다.
    \item \textbf{파산:} 리먼 브라더스는 2008년 금융 위기 동안 파산했습니다. 이 소식이 전해지자 공황에 빠진 투자자들은 펀드에서 돈을 인출하기 위해 몰려들었습니다.
    \item \textbf{버크 깨짐 (Breaking the Buck):} 결국 펀드는 투자 주당 1달러로 모든 상환을 충족시킬 수 없었습니다. 파산에 들어갔습니다.
    \textbf{결과:} 리저브 프라이머리 펀드는 결국 주당 1달러당 0.97달러만 돌려줄 수 있었습니다. 이 사건을 \textbf{버크 깨짐 (breaking the buck)}이라고 합니다. 머니 마켓 뮤추얼 펀드에는 예금 보험이 없으므로, 펀드가 버크를 깨면 투자자들은 돈을 잃게 됩니다.
    \item \textbf{금융 시장에 미친 연쇄 반응:}
    \begin{itemize}
        \item 리저브 프라이머리 펀드가 실패한 후, 투자자들은 다른 프라임 펀드의 건전성에 대해 우려하게 되었습니다. 안전해 보이던 투자가 갑자기 위험해 보였습니다.
        \item 그들은 이 펀드에서 약 4천억 달러를 인출했습니다. 이는 머니 마켓의 현금 감소로 이어졌습니다.
        \item 머니 마켓 펀드는 단기 자금 시장에서 철수할 수밖에 없었습니다.
        \item 이 시장에서 보통 돈을 빌리던 기업과 금융 기관들은 현금을 확보하거나 자산을 매각하기 위해 분주했습니다. 은행의 경우, 기업과 가계에 대한 신용 공급을 줄였습니다.
    \end{itemize}
\end{itemize}

\begin{summarybox}[title=\textbf{리저브 프라이머리 펀드 사례의 교훈}]
리저브 프라이머리 펀드의 실패는 머니 마켓 펀드가 금융 시스템에 얼마나 중요한지를 보여줍니다. 이 펀드들이 어려움을 겪고 대규모 상환이 발생하면, 이는 금융 시스템의 많은 부분에 매우 문제가 될 수 있습니다. 금융 위기 이후, 금융 시장 규제 당국은 머니 마켓 뮤추얼 펀드의 미래 혼란을 피하기 위한 개혁을 시행했습니다.
\end{summarybox}

\subsubsection*{핵심 요약: 머니 마켓 뮤추얼 펀드}
\begin{itemize}
    \item 머니 마켓 뮤추얼 펀드는 미국 금융 시스템에서 가장 큰 단기 또는 현금 대출 기관입니다. 이 펀드들은 경제 전반에 걸쳐 현금을 효율적으로 배분하는 데 도움을 줍니다.
    \item 투자자들은 머니 마켓 뮤추얼 펀드를 은행 예금의 대체재로 봅니다. 그러나 이러한 투자는 은행 예금처럼 보험에 가입되어 있지 않습니다.
    \item 이 펀드들의 혼란은 단기 신용 감소를 통해 금융 시스템의 많은 부분에 영향을 미칠 수 있습니다.
\end{itemize}

\subsection{증권화 (Securitization)}
\begin{warningbox}[title=\textbf{참고: 증권화 (Securitization) 내용}]
이 문서의 파트 1/2에는 'Lesson 1-3.2. Securitization'이 목차에 포함되어 있으나, 해당 강의의 본문 내용은 제공되지 않았습니다. 따라서 이 섹션에는 추가적인 설명이 포함되어 있지 않습니다.
\end{warningbox}

\newpage
\section{빠르게 훑어보기 (1페이지 요약)}

\begin{summarybox}[title=\textbf{모듈 1 핵심 요약: 금리 및 금융 시장}]
\begin{itemize}
    \item \textbf{중앙은행의 역할:} 통화정책은 금리 조정을 통해 경제에 영향을 미치며, 그 핵심은 단기 금리 시장에 있습니다.
    \item \textbf{단기 금리 시장:}
    \begin{itemize}
        \item \textbf{연방기금 시장:} 은행 간 초과 지급준비금 거래 시장. 연방기금 금리는 중앙은행의 주요 정책 목표 금리이자 모든 금리에 영향을 미치는 기준점입니다.
        \item \textbf{환매조건부채권 (Repo) 시장:} 담보(주로 국채)를 이용한 익일물 대출 시장. 모든 금융 기관이 참여하며, 헤어컷을 통해 대출자의 위험을 줄입니다.
        \item \textbf{LIBOR vs. SOFR:} LIBOR는 무담보 은행 간 금리로 조작 논란 후 SOFR(담보부 익일물 자금 조달 금리)로 대체되고 있습니다. SOFR은 실제 거래 기반의 무위험 금리입니다.
    \end{itemize}
    \item \textbf{장기 금리 개념:}
    \begin{itemize}
        \item \textbf{명목 금리 vs. 실질 금리:} 명목 금리는 인플레이션을 고려하지 않은 표면적 금리이며, 실질 금리는 인플레이션을 조정한 실제 구매력 변화를 반영합니다 (피셔 방정식: 명목 금리 $\approx$ 실질 금리 + 예상 인플레이션).
        \item \textbf{물가연동국채 (TIPS):} 원금이 인플레이션에 따라 조정되어 실질 금리를 보장하며, 시장의 미래 인플레이션 기대를 측정하는 중요한 지표입니다.
    \end{itemize}
    \item \textbf{국채와 무위험 금리:}
    \begin{itemize}
        \item \textbf{미국 국채:} 세계 최대, 가장 유동적인 채권 시장. 신용 위험이 없어 무위험 금리의 벤치마크로 사용됩니다. T-Bill, T-Note, T-Bond로 구분됩니다.
        \item \textbf{국채 투자 위험:} 인플레이션과 금리 변화 위험이 있으며, 장기 국채의 명목 금리에는 실질 금리, 예상 인플레이션, 기간 프리미엄이 포함됩니다.
    \end{itemize}
    \item \textbf{금리 기간 구조 (수익률 곡선):}
    \begin{itemize}
        \item 만기별 금리를 비교한 것으로, 경제 상황과 미래 금리 기대를 반영합니다.
        \item \textbf{우상향 곡선:} 경제 확장 및 금리 인상 기대.
        \item \textbf{역전된 곡선:} 경기 침체 및 금리 인하 기대의 강력한 예측 지표.
    \end{itemize}
    \item \textbf{기타 자금 시장:}
    \begin{itemize}
        \item \textbf{머니 마켓 뮤추얼 펀드:} 초안전 단기 자산에 투자하는 펀드로, 은행 예금의 대안이자 단기 자금 시장의 핵심입니다. 정부 보험이 없어 '버크 깨짐' 위험이 있습니다 (예: 리저브 프라이머리 펀드 사례).
        \item \textbf{증권화:} (본 문서에서는 내용 미제공)
    \end{itemize}
\end{itemize}
\end{summarybox}

\newpage

\newpage
\section{개요: 증권화 (Securitization)}

\begin{tcolorbox}[
    colback=lightblue,
    colframe=darkblue,
    title=\textbf{문서 핵심 요약}
]
  이 문서는 금융 혁신인 \textbf{증권화 (Securitization)}에 대해 다룹니다.
  증권화는 1970년대에 부상하여 미국 경제에 지대한 영향을 미쳤습니다.
  수많은 대출을 한데 모아 \textbf{위험을 분산}하고 \textbf{현금 흐름을 정확하게 예측}함으로써,
  이 대출들을 담보로 \textbf{증권(채권과 지분)}을 발행하여 자금을 조달하는 방식입니다.
  이를 통해 대출 공급이 확대되었지만, \textbf{단기 자금 시장의 변동성}에 민감하다는 단점도 있습니다.
\end{tcolorbox}

\newpage
\section{용어 정리 (증권화)}

\begin{tcolorbox}[
    colback=lightgreen,
    colframe=darkgreen,
    title=\textbf{핵심 용어 해설}
]
  다음 표는 증권화 개념을 이해하는 데 필수적인 용어들을 정리한 것입니다.
  각 용어의 쉬운 설명과 원어를 함께 제시하여 이해를 돕습니다.

  \begin{adjustbox}{width=\textwidth,center}
  \begin{tabular}{lllc}
  \toprule
  \textbf{용어 (한글)} & \textbf{쉬운 설명} & \textbf{원어} & \textbf{비고} \\
  \midrule
  \textbf{증권화} & 대출과 같은 비유동성 자산을 모아 유동성 있는 증권(채권, 주식)으로 전환하여 투자자에게 판매하는 과정. & Securitization & 금융 혁신의 핵심 \\
  \textbf{모기지 대출} & 부동산을 담보로 제공하고 빌리는 대출. & Mortgage Loan & 주택 구매 시 주로 이용 \\
  \textbf{채무 불이행} & 대출금을 약속대로 상환하지 못하는 것. & Default & 대출의 주요 위험 요소 \\
  \textbf{대출 풀} & 여러 개의 대출을 한데 모아 놓은 집합. & Loan Pool & 위험 분산의 기본 단위 \\
  \textbf{다각화} & 여러 자산에 분산 투자하여 특정 자산의 위험이 전체 포트폴리오에 미치는 영향을 줄이는 전략. & Diversification & '모든 달걀을 한 바구니에 담지 마라' \\
  \textbf{현금 흐름} & 일정 기간 동안 기업이나 자산에서 발생하고 유출되는 현금의 총량. & Cash Flow & 수익성과 안정성 판단 기준 \\
  \textbf{증권} & 주식, 채권 등 재산적 가치를 나타내는 문서 또는 전자 기록. & Securities & 투자 및 자금 조달 수단 \\
  \textbf{채권} & 정부, 기업 등이 자금을 조달하기 위해 발행하는 차용증서. 정해진 이자를 지급하고 만기에 원금을 상환. & Bonds & 비교적 안전한 투자 \\
  \textbf{지분} & 기업의 소유권을 나타내는 부분. 채권자에게 지급하고 남은 이익에 대한 청구권을 가짐. & Equity & 고위험 고수익 가능성 \\
  \textbf{쿠폰 지급} & 채권 투자자에게 정기적으로 지급되는 이자. & Coupon Payments & 채권 수익의 한 형태 \\
  \textbf{우선 청구권} & 특정 자산이나 현금 흐름에 대해 다른 채권자보다 먼저 권리를 행사할 수 있는 권리. & First Claim/Priority & 채권자가 지분보다 우선 \\
  \textbf{기업어음} & 기업이 단기 자금을 조달하기 위해 발행하는 무담보 어음. & Commercial Paper (CP) & 만기가 짧은 단기 채권 \\
  \textbf{머니마켓 뮤추얼 펀드} & 단기 금융 상품에 투자하여 수익을 추구하는 펀드. & Money Market Mutual Funds (MMMF) & 단기 자금 시장의 주요 참여자 \\
  \textbf{정부 후원 기업} & 정부가 설립하거나 후원하여 특정 공공 목적(예: 주택 금융)을 달성하는 데 도움을 주는 민간 기업. & Government Sponsored Enterprises (GSEs) & Fannie Mae, Freddie Mac 등 \\
  \textbf{미국 국채} & 미국 정부가 발행하는 채권. 가장 안전한 투자 자산 중 하나로 간주됨. & US Treasury Securities & 무위험 자산의 대명사 \\
  \bottomrule
  \end{tabular}
  \end{adjustbox}
\end{tcolorbox}

\newpage
\section{핵심 개념 및 원리: 증권화의 이해}

\subsection{1. 증권화란 무엇인가?}

증권화는 1970년대에 등장하여 금융 시장에 큰 변화를 가져온 중요한 금융 혁신입니다. 이는 단순히 대출을 묶는 것을 넘어, 대출의 위험과 수익을 분리하고 재구성하여 새로운 투자 상품을 만들어내는 복잡한 과정입니다.

\begin{tcolorbox}[
    colback=lightyellow,
    colframe=darkorange,
    title=\textbf{증권화의 정의와 중요성}
]
  \textbf{증권화 (Securitization)}는 은행이나 다른 금융 기관이 보유한 대출(예: 모기지, 자동차 대출, 학자금 대출 등)과 같은 \textbf{비유동성 자산}들을 한데 모아(풀, Pool), 이 자산들에서 발생하는 미래의 현금 흐름(원금 및 이자 상환)을 담보로 \textbf{유동성 있는 증권(Securities)}을 발행하여 투자자들에게 판매하는 과정입니다.

  이러한 증권화는 대출 기관이 대출을 계속해서 제공할 수 있도록 새로운 자금 조달원을 제공하며, 투자자들에게는 다양한 위험-수익 프로필을 가진 새로운 투자 기회를 제공합니다. 특히 미국 경제에서 신용 공급을 확대하는 데 매우 중요한 역할을 해왔습니다.
\end{tcolorbox}

\subsection{2. 증권화의 작동 원리: 불확실성 감소}

증권화의 핵심은 \textbf{대규모의 대출을 모아 위험을 분산}하고, 이를 통해 미래의 현금 흐름을 \textbf{더욱 정확하게 예측}할 수 있게 만드는 것입니다. 이 원리를 이해하기 위해 대출의 규모에 따른 위험 변화를 살펴보겠습니다.

\subsubsection{2.1. 단일 대출의 높은 불확실성}

하나의 고객에게 하나의 모기지 대출을 제공했다고 가정해 봅시다. 고객의 신용도(위험 프로필)를 분석한 결과, 이 고객이 채무 불이행(Default)할 확률이 2\%라고 예상됩니다.

\begin{tcolorbox}[
    colback=boxred,
    colframe=red!70!black,
    title=\textbf{단일 대출의 위험}
]
  단 하나의 대출과 하나의 고객만 있다면, 결과는 둘 중 하나입니다.
  \begin{itemize}
    \item 고객이 대출을 \textbf{완전히 상환}하거나,
    \item 고객이 대출을 \textbf{완전히 불이행}하거나.
  \end{itemize}
  이처럼 단일 대출의 손실은 \textbf{매우 불확실}합니다. '모 아니면 도' (all or nothing)의 상황이기 때문에, 대출 기관은 큰 위험에 직면하게 됩니다. 2\%의 불이행 확률은 낮아 보이지만, 실제로 불이행이 발생하면 전체 대출금을 잃게 됩니다.
\end{tcolorbox}

\subsubsection{2.2. 대출 풀(Pool)을 통한 위험 분산}

이제 이러한 모기지 대출 100개를 모아 '대출 풀(Loan Pool)'을 만들었다고 가정해 봅시다. 각 고객의 채무 불이행 확률은 여전히 2\%입니다. 중요한 것은 이 대출자들이 모두 다른 지역에 살고 있다고 가정하는 것입니다. 이는 한 지역의 경제적 충격이 모든 대출자에게 동시에 영향을 미치지 않도록 하여, \textbf{동시에 채무 불이행할 가능성이 낮다}는 것을 의미합니다.

\begin{tcolorbox}[
    colback=lightblue,
    colframe=darkblue,
    title=\textbf{대출 풀의 위험 감소 효과}
]
  100개의 대출 풀에서 채무 불이행 위험은 어떻게 될까요?
  \begin{itemize}
    \item 평균적으로 100명 중 2명(2\%)이 대출을 상환하지 않을 것으로 예상합니다.
    \item 하지만 실제로는 1명이 불이행할 수도 있고, 3명이 불이행할 수도 있습니다.
    \item 중요한 점은 \textbf{'모 아니면 도'가 아니라는 것}입니다. 100개의 대출 중 일부가 불이행하더라도, 나머지 대출들은 정상적으로 상환됩니다.
  \end{itemize}
  이처럼 100개의 모기지 대출 풀은 단일 대출 예시보다 \textbf{불확실성이 훨씬 줄어듭니다}.

  더 나아가, 10,000개의 모기지 대출 풀을 생각해 봅시다. 각 대출의 불이행 확률은 여전히 2\%이므로, 평균적으로 200명(2\%)이 불이행할 것으로 예상됩니다. 이 경우, 실제 불이행 건수는 190건 또는 210건이 될 수 있으며, 이는 전체의 1.9\% 또는 2.1\%에 해당합니다. 이 수치는 우리가 예상한 2\%에 \textbf{매우 근접}합니다.
\end{tcolorbox}

\subsubsection{2.3. 핵심 원리: 대규모 분산 풀의 예측 가능성}

위의 예시에서 얻을 수 있는 핵심 교훈은 다음과 같습니다.

\begin{tcolorbox}[
    colback=lightgreen,
    colframe=darkgreen,
    title=\textbf{증권화의 핵심: 현금 흐름의 정확한 예측}
]
  우리는 \textbf{크고 잘 다각화된(well-diversified) 모기지 대출 풀}을 원합니다.
  \begin{itemize}
    \item '크다'는 것은 대출의 수가 많다는 의미입니다.
    \item '잘 다각화되었다'는 것은 대출자들이 서로 다른 지역에 살거나, 다른 산업에 종사하는 등, 특정 요인에 의해 동시에 영향을 받을 가능성이 낮다는 의미입니다.
  \end{itemize}
  이러한 대규모의 잘 다각화된 풀을 통해 우리는 대출에서 발생하는 \textbf{현금 흐름(Cash Flow)을 매우 정확하게 추정}할 수 있습니다. 현금 흐름은 대출의 손실(불이행)과 상환액을 모두 포함합니다.

  대출 풀의 현금 흐름을 정확하게 예측할 수 있다는 점이 바로 \textbf{증권화가 작동하는 핵심 원리}입니다. 예측 가능성이 높을수록 투자자들은 더 안심하고 투자할 수 있게 됩니다.
\end{tcolorbox}

\subsection{3. 증권 발행: 채권과 지분}

이제 증권화의 두 번째 중요한 요소인 \textbf{증권 발행}에 대해 알아보겠습니다. 증권화 과정에서 발행 기관(Issuer)은 대출 풀을 담보로 증권을 판매합니다. 이 증권은 주로 \textbf{채권(Bonds)}과 \textbf{지분(Equity)}으로 구성됩니다. 이 증권들을 판매하여 얻은 자금은 모기지 대출에 필요한 자금을 조달하는 데 사용됩니다.

\begin{tcolorbox}[
    colback=lightyellow,
    colframe=darkorange,
    title=\textbf{증권의 구성과 우선순위}
]
  \begin{itemize}
    \item \textbf{채권 (Bonds):}
      \begin{itemize}
        \item 대출 풀의 현금 흐름을 정확하게 추정할 수 있다면, 상당한 양의 채권을 발행할 수 있습니다.
        \item 이 채권에 대한 이자(쿠폰 지급, Coupon Payments)는 대출 풀에서 발생하는 모기지 이자 상환액에서 나옵니다.
        \item 채권은 \textbf{가장 먼저 현금 흐름에 대한 청구권(First Claim)}을 가집니다. 즉, 대출 풀에서 돈이 들어오면 채권 보유자들에게 먼저 지급됩니다.
        \item 이러한 우선순위 때문에 채권은 \textbf{상대적으로 안전한 투자 상품}으로 간주됩니다.
      \end{itemize}
    \item \textbf{지분 (Equity):}
      \begin{itemize}
        \item 지분 증권 또한 모기지 풀에서 현금 흐름을 받습니다.
        \item 하지만 지분 보유자들은 \textbf{두 번째로 지급}받습니다. 모든 채권 보유자들에게 지급이 완료된 후에 남은 현금 흐름을 지분 보유자들이 가져갑니다.
        \item 이 때문에 지분은 채권보다 \textbf{더 높은 위험}을 가지지만, 동시에 \textbf{더 높은 수익}을 얻을 가능성도 있습니다.
      \end{itemize}
  \end{itemize}
  이러한 구조는 다양한 위험 선호도를 가진 투자자들을 유치할 수 있게 합니다. 안전한 투자를 선호하는 투자자는 채권을, 높은 수익을 추구하는 투자자는 지분을 선택할 수 있습니다.
\end{tcolorbox}

\subsection{4. 증권화의 수익성: 구체적인 예시}

증권화가 왜 수익성이 높은지 구체적인 예를 통해 살펴보겠습니다.

\begin{tcolorbox}[
    colback=lightblue,
    colframe=darkblue,
    title=\textbf{증권화 수익성 예시}
]
  \textbf{가정:}
  \begin{itemize}
    \item 100달러 가치의 모기지 대출 풀이 있습니다.
    \item 이 모기지 대출들은 연 5\%의 이자를 지급합니다.
    \item 우리는 이 풀을 담보로 80달러의 채권을 발행하고, 20달러의 지분을 직접 보유합니다.
    \item 발행된 총 100달러의 증권(채권 80달러 + 지분 20달러)이 100달러의 모기지 대출에 자금을 조달합니다.
  \end{itemize}

  채권 보유자들은 우선적으로 지급받기 때문에 상대적으로 안전합니다. 우리는 채권에 연 4\%의 이자를 약속했다고 가정합시다.

  \textbf{시나리오 1: 채무 불이행이 전혀 없는 경우}
  \begin{itemize}
    \item 모기지 풀은 총 5달러의 현금(100달러의 5\%)을 발생시킵니다.
    \item 80달러의 채권은 4\%의 수익률을 받으므로, 채권 보유자들은 3.20달러 (80달러 $\times$ 4\%)를 받습니다.
    \item 지분 보유자들은 남은 1.80달러 (5달러 - 3.20달러)를 받습니다.
  \end{itemize}
  이 경우, 지분 투자액 20달러 대비 1.80달러의 수익은 9\% (1.80/20)로, 채권 수익률 4\%보다 훨씬 높습니다.

  \textbf{시나리오 2: 모기지 중 2\%가 채무 불이행하는 경우}
  \begin{itemize}
    \item 모기지 풀은 총 4.90달러의 현금 (100달러 중 2달러가 불이행하여 98달러만 상환된다고 가정하면, 98달러의 5%는 4.90달러)을 발생시킵니다.
    \item 채권 보유자들은 우선순위가 있으므로 여전히 3.20달러를 받습니다.
    \item 지분 보유자들은 남은 1.70달러 (4.90달러 - 3.20달러)를 받습니다. 이는 시나리오 1보다 0.10달러 적은 금액입니다.
  \end{itemize}

  \textbf{시나리오 3: 예상치 못한 높은 채무 불이행이 발생하는 경우}
  \begin{itemize}
    \item 만약 채무 불이행이 매우 높아서 모기지 풀에서 발생하는 현금 흐름이 3.20달러 이하로 떨어진다면, 채권 보유자들은 약속된 이자를 받지 못할 수도 있습니다.
    \item 하지만 그 전에, 지분 보유자들은 아무것도 받지 못할 수 있습니다. 지분 보유자들은 일반적으로 투자에 대해 더 높은 수익을 기대하지만, 그만큼 더 큰 위험을 감수합니다.
  \end{itemize}
  이 예시는 증권화의 기본 원리를 보여줍니다. 대출 풀의 현금 흐름을 정확하게 추정할 수 있다면, 채권을 사용하여 대출에 자금을 조달할 수 있습니다. 대출 기관은 모든 대출에 필요한 현금을 직접 보유할 필요 없이 시장에서 자금을 빌릴 수 있게 됩니다.
\end{tcolorbox}

\newpage
\section{증권화의 중요성과 영향}

\subsection{1. 자금 조달 모델로서의 증권화}

증권화는 20세기 가장 중요한 금융 혁신 중 하나로 평가받습니다. 이는 단순한 투자 상품을 넘어, 금융 기관의 자금 조달 방식 자체를 변화시켰습니다.

\begin{tcolorbox}[
    colback=lightgreen,
    colframe=darkgreen,
    title=\textbf{증권화가 신용 공급에 미치는 영향}
]
  \begin{itemize}
    \item \textbf{대출 기관의 부담 경감:} 전통적인 방식에서는 은행이 대출을 제공하려면 그만큼의 예금이나 자체 자본을 보유해야 했습니다. 하지만 증권화를 통해 대출 기관은 대출을 제공한 후 이를 증권화하여 시장에 판매함으로써, 필요한 자금을 투자자들로부터 조달할 수 있게 됩니다.
    \item \textbf{신용 공급 확대:} 대출 기관이 자체 자본에 묶이지 않고 지속적으로 자금을 조달할 수 있게 되면서, 기업과 가계에 대한 대출 공급 능력이 크게 증가했습니다. 1970년대 이후 미국에서 기업과 가계에 대한 신용 공급이 증권화 시장 덕분에 크게 늘어났습니다.
    \item \textbf{자본 효율성 증대:} 은행은 대출을 증권화하여 장부에서 제거함으로써, 대출에 묶여 있던 자본을 다른 대출이나 투자에 활용할 수 있게 되어 자본 효율성을 높일 수 있습니다.
  \end{itemize}
  결론적으로, 증권화는 금융 시스템 내에서 자금이 더욱 효율적으로 순환하고, 더 많은 사람들이 대출을 받을 수 있도록 하는 중요한 메커니즘이 되었습니다.
\end{tcolorbox}

\subsection{2. 미국 내 증권화 현황 데이터}

미국에서는 가계 및 기업 대출의 상당 부분이 증권화되어 있습니다. 이 현상은 1980년대 이후 극적으로 증가했습니다.

\begin{tcolorbox}[
    colback=lightblue,
    colframe=darkblue,
    title=\textbf{미국 증권화 시장의 규모와 특징}
]
  \begin{itemize}
    \item \textbf{시장 규모:} 2008년 금융 위기 이후 잠시 감소했지만, 2021년까지 증권화 관련 채권의 총액은 약 \textbf{12조 달러}에 달했습니다. 이는 엄청난 규모의 금융 시장을 형성하고 있음을 보여줍니다.
    \item \textbf{주요 발행 주체:}
      \begin{itemize}
        \item 이 채권의 대부분, 즉 \textbf{10조 달러 이상}은 \textbf{정부 후원 기업 (Government Sponsored Enterprises, GSEs)}에 의해 발행되었으며, 미국 정부의 보증을 받습니다.
        \item 정부 보증 덕분에 이 채권들은 \textbf{미국 국채 (US Treasury Securities)만큼 안전}하다고 평가됩니다.
        \item 나머지 부분은 은행 및 금융 회사와 같은 \textbf{민간 발행 기관}의 증권화에서 비롯됩니다.
      \end{itemize}
    \item \textbf{주요 증권화 대상:}
      \begin{itemize}
        \item 증권화된 대출의 \textbf{대부분은 모기지}입니다. 여기에는 단독 주택을 위한 주거용 모기지, 다세대 주택 및 아파트 건물을 위한 모기지, 사무실 건물이나 쇼핑몰을 위한 상업용 모기지가 포함됩니다.
        \item 다른 대출 범주에는 \textbf{자동차 대출}과 \textbf{소기업 대출} 등이 있습니다.
      \end{itemize}
  \end{itemize}
  이 데이터는 증권화가 미국 금융 시스템의 핵심적인 부분이며, 특히 주택 금융 시장에서 절대적인 영향력을 가지고 있음을 명확히 보여줍니다.
\end{tcolorbox}

\subsection{3. 단기 자금 시장과의 민감한 관계}

증권화는 매우 효율적인 자금 조달 모델이지만, \textbf{단기 자금 시장 (Short-term Funding Markets)}의 상황에 민감하게 반응할 수 있다는 중요한 단점이 있습니다.

\begin{tcolorbox}[
    colback=boxred,
    colframe=red!70!black,
    title=\textbf{단기 자금 시장과 증권화의 연결고리}
]
  \textbf{왜 증권화가 단기 자금 시장에 민감할까요?}
  \begin{itemize}
    \item \textbf{임시 자금 조달:} 많은 발행 기관들은 대출을 풀로 모아 증권화하기 전에, \textbf{단기적으로 자금을 조달}하여 대출을 제공합니다.
    \item \textbf{기업어음 (Commercial Paper) 활용:} 예를 들어, 이들은 머니마켓 뮤추얼 펀드(Money Market Mutual Funds)에 기업어음(Commercial Paper)을 판매하여 현금을 조달합니다. 이 현금으로 대출을 제공합니다.
    \item \textbf{상환 과정:} 충분한 대출이 모여 풀이 형성되고 증권화되면, 발행 기관은 채권 투자자들로부터 조달한 자금으로 이전에 발행했던 기업어음을 상환합니다.
  \end{itemize}
  이러한 과정은 단기 자금 시장이 원활하게 작동할 때만 가능합니다.

  \textbf{단기 자금 시장의 교란이 경제에 미치는 영향:}
  \begin{itemize}
    \item 만약 단기 자금 시장이 \textbf{경색(tries up)}되면, 즉 기업어음과 같은 단기 자금 조달이 어려워지면, 증권화 모델은 작동을 멈추게 됩니다.
    \item 2008년 금융 위기 때처럼 이러한 상황이 발생하면, 기업과 가계는 대출을 받기 어려워져 \textbf{신용 경색 (Credit Crunch)}이 발생하고 경제 전반에 심각한 영향을 미칠 수 있습니다.
  \end{itemize}
  따라서 증권화는 신용 공급을 확대하는 강력한 도구이지만, 동시에 금융 시스템의 취약성을 증가시킬 수 있는 양날의 검과 같습니다.
\end{tcolorbox}

\newpage
\section{세션 요약 (증권화)}

이번 세션에서는 금융 시장의 중요한 혁신인 증권화에 대해 깊이 있게 학습했습니다. 핵심 내용은 다음과 같습니다.

\begin{tcolorbox}[
    colback=lightgreen,
    colframe=darkgreen,
    title=\textbf{학습 내용 요약}
]
  \begin{itemize}
    \item \textbf{증권화의 정의와 작동 원리:}
      \begin{itemize}
        \item 증권화는 20세기에 등장한 중요한 자금 조달 모델입니다.
        \item 이는 \textbf{통계 모델}에 기반하며, 많은 수의 대출을 모아 \textbf{불확실성을 줄이는 것}이 핵심입니다.
        \item 대출 풀을 형성하고, 이 풀을 담보로 \textbf{증권(채권과 지분)}을 발행하는 과정을 포함합니다.
      \end{itemize}
    \item \textbf{증권화의 긍정적 영향:}
      \begin{itemize}
        \item 증권화 자금 조달 모델은 미국에서 \textbf{총 신용 공급을 증가}시키는 데 크게 기여했습니다.
        \item 이를 통해 기업과 가계가 더 쉽게 대출을 받을 수 있게 되었습니다.
      \end{itemize}
    \item \textbf{증권화의 잠재적 단점:}
      \begin{itemize}
        \item 증권화는 \textbf{단기 자금 시장의 상황에 민감}할 수 있습니다.
        \item 단기 자금 시장의 교란은 증권화 모델의 붕괴로 이어져, 신용 경색을 유발하고 경제에 부정적인 영향을 미칠 수 있습니다.
      \end{itemize}
  \end{itemize}
  증권화는 금융 시스템의 효율성을 높이는 동시에, 잠재적인 위험 요소를 내포하고 있음을 이해하는 것이 중요합니다.
\end{tcolorbox}

\end{document}
